% !TeX program = pdflatex
% !BIB program = bibtex
% =============================================================================
% Manuscript on the paper template of the Chair of Computational Mechanics,
% OVGU Magdeburg. Compiles with pdfLaTeX and LuaLaTeX:
%   latexmk -pdf main        (or: latexmk -lualatex main)
% =============================================================================
\documentclass{article}
%
% -----------------------------------------------------------------------------
% Layout: page geometry, header, title block (do not change)
% -----------------------------------------------------------------------------
\usepackage{arxiv}
%
% -----------------------------------------------------------------------------
% Paper metadata: edit here only (used for title, header and PDF metadata)
% -----------------------------------------------------------------------------
\newcommand{\papertitle}{Unsymmetric spectral elements: Stable metric shape functions based on a hybrid metric-parametric trial space}
\newcommand{\papershorttitle}{}  % centre of the page header; empty, since the long
                                % journal name on the right extends past the centre
\newcommand{\journalname}{Computer Methods in Applied Mechanics and Engineering}  % target journal
\newcommand{\paperauthors}{S.~Eisentr\"ager}  % PDF metadata only
\newcommand{\paperkeywords}{unsymmetric finite element method, Petrov-Galerkin method, spectral element method, curved-edge distortion, metric shape functions, Lebesgue constant, generalized Vandermonde matrix}
%
% -----------------------------------------------------------------------------
% Fonts and encoding (UTF-8 input is the LaTeX default since 2018)
% Times-like text and math (newtx), Inconsolata as typewriter font; identical
% output with pdfLaTeX (Type 1 fonts) and LuaLaTeX (OpenType via fontspec).
% -----------------------------------------------------------------------------
\usepackage{iftex}
\ifPDFTeX
  \usepackage[T1]{fontenc}
\fi
\usepackage{amsmath}    % must be loaded BEFORE newtxmath
\usepackage{amsthm}     % theorem environments (defined in preamble.tex)
\usepackage{newtxtext}  % with LuaLaTeX, this loads fontspec
\usepackage{newtxmath}  % includes the AMS symbols: do NOT load amssymb/amsfonts
\ifPDFTeX
  \usepackage{inconsolata}  % after newtxtext, which sets its own \ttdefault
\else
  \setmonofont{Inconsolatazi4-Regular.otf}[BoldFont=Inconsolatazi4-Bold.otf]
  \setsansfont{texgyreheros-regular.otf}[
    BoldFont=texgyreheros-bold.otf, ItalicFont=texgyreheros-italic.otf,
    BoldItalicFont=texgyreheros-bolditalic.otf]  % as newtxtext with pdfLaTeX
  \DeclareMathAlphabet{\mathsf}{\encodingdefault}{\sfdefault}{m}{n}
\fi
\usepackage{microtype}
%
% -----------------------------------------------------------------------------
% Mathematics
% -----------------------------------------------------------------------------
\usepackage{mathtools}
\usepackage{bm}         % bold math symbols: \bm{\sigma}
\usepackage{cancel}
\usepackage{xfrac}      % compact fractions: \sfrac{1}{2}
\usepackage{scalerel}
\usepackage{mleftright} % \left( ... \right) without extra space
\mleftright             % ... applied to all \left/\right pairs
%
% -----------------------------------------------------------------------------
% Tables (booktabs rules: \toprule, \midrule, \cmidrule, \bottomrule)
% -----------------------------------------------------------------------------
\usepackage{booktabs}
\usepackage{longtable}
\usepackage{colortbl}
\usepackage{arydshln}   % dashed lines; load after longtable and colortbl
\usepackage{diagbox}
\usepackage{multirow}
\usepackage{adjustbox}  % wide tables scaled to the text width
\usepackage{siunitx}    % units \qty{210}{\giga\pascal}, number columns S
\sisetup{
  exponent-product     = \cdot,   % 1.2 \cdot 10^{-3} instead of 1.2 \times 10^{-3}
  separate-uncertainty = true,
  range-phrase         = --,
}
%
% -----------------------------------------------------------------------------
% Figures
% -----------------------------------------------------------------------------
\usepackage{graphicx}
\graphicspath{{./figs/}}
\usepackage{overpic}    % LaTeX labels on top of graphics
\usepackage{caption}
\captionsetup[table]{position=top, skip=7pt}  % caption above tables
\usepackage{subcaption} % subfigures: subfigure environment or \subfloat
\usepackage{pgfplots}   % plots from data files
\pgfplotsset{compat=1.18}
\usepgfplotslibrary{groupplots}
\usepackage{placeins}   % \FloatBarrier
%
% -----------------------------------------------------------------------------
% Algorithms
% -----------------------------------------------------------------------------
\usepackage{algorithm}
\usepackage{algpseudocode}
%
% -----------------------------------------------------------------------------
% Footnotes, citations, line numbers
% -----------------------------------------------------------------------------
\usepackage[bottom,symbol]{footmisc}  % footnotes below floats
\usepackage[noadjust]{cite}           % sorted and compressed [1-3]
\usepackage{lineno}                   % \linenumbers (see below)
%
% -----------------------------------------------------------------------------
% Hyperlinks (load late) and clever references (load last)
% -----------------------------------------------------------------------------
\usepackage{hyperref}
\hypersetup{
  pdftitle    = {\papertitle},
  pdfauthor   = {\paperauthors},
  pdfsubject  = {Computational Mechanics},
  pdfkeywords = {\paperkeywords},
  hidelinks,
}
\usepackage[capitalise]{cleveref}  % \cref: Fig. 1, Eq. (1), Table 1, Section 1;
                                    % \Cref at the start of a sentence: Figure 1, Equation (1)
%
% -----------------------------------------------------------------------------
% Notation: define symbols once, use them consistently
% -----------------------------------------------------------------------------
\newcommand{\vect}[1]{\bm{#1}}                       % vectors and matrices
\newcommand{\dd}{\mathop{}\!\mathrm{d}}              % differential: \dd x
\DeclareRobustCommand{\transp}{\mathsf{T}}             % transpose: \vect{u}^\transp (overrides newtxmath)
\newcommand{\dof}{\mathrm{DOF}}
%
% Paper-specific theorem environments, notation and plot styles
\usepackage[T1]{fontenc}
\usepackage[utf8]{inputenc}
\usepackage{lmodern}
\usepackage{amsmath,amssymb,amsthm}
\usepackage{bm}
\usepackage{booktabs}
\usepackage{multirow}
\usepackage{adjustbox}
\usepackage{graphicx}
\usepackage{xcolor}
\usepackage{tikz}
\usepackage{pgfplots}
\usepgfplotslibrary{groupplots}
\pgfplotsset{compat=1.18}
\usepackage{algorithm}
\usepackage{algpseudocode}
\usepackage[colorlinks=true,linkcolor=blue!50!black,citecolor=blue!50!black,
  urlcolor=blue!50!black]{hyperref}
\pdfstringdefDisableCommands{\def\corref#1{}}  % no footnote mark in the PDF metadata
\usepackage[capitalise]{cleveref}
\crefname{equation}{Eq.}{Eqs.}\Crefname{equation}{Equation}{Equations}
\crefname{section}{Sect.}{Sects.}\Crefname{section}{Section}{Sections}
\crefname{subsection}{Sect.}{Sects.}\Crefname{subsection}{Section}{Sections}
\crefname{figure}{Fig.}{Figs.}\Crefname{figure}{Figure}{Figures}
\crefname{table}{Table}{Tables}\Crefname{table}{Table}{Tables}
\crefname{appendix}{Appendix}{Appendices}\Crefname{appendix}{Appendix}{Appendices}
\crefname{algorithm}{Algorithm}{Algorithms}

% Theorem environments
\newtheorem{proposition}{Proposition}
\newtheorem{lemma}{Lemma}
\newtheorem{corollary}{Corollary}
\theoremstyle{definition}
\newtheorem{definition}{Definition}
\theoremstyle{remark}
\newtheorem{remark}{Remark}

% Notation
\newcommand{\vect}[1]{\bm{#1}}
\newcommand{\matr}[1]{\bm{#1}}
\newcommand{\T}{\top}
\newcommand{\dd}{\,\mathrm{d}}
\newcommand{\bx}{\vect x}
\newcommand{\bxi}{\bm\xi}
\newcommand{\R}{\mathbb{R}}
\newcommand{\Pp}{\mathcal P}
\newcommand{\Qp}{\mathcal Q}
\newcommand{\Tp}{\mathcal T}
\newcommand{\Cp}{\mathcal C}
\newcommand{\Ihat}{\hat\Pi}
\newcommand{\Lam}{\Lambda}
\newcommand{\adj}{\operatorname{adj}}
\newcommand{\spn}{\operatorname{span}}
\newcommand{\cond}{\kappa}
% Notation of the companion paper (Eisentraeger et al., axisymmetric solids)
\newcommand{\Kh}{\hat K}
\newcommand{\bxib}{\bar{\bm\xi}}
\newcommand{\bxit}{\tilde{\bm\xi}}
\newcommand{\Qsk}{\Qp_p(\bar\xi,\bar\eta)}
\newcommand{\nG}{n_{\mathrm{G}}}
\newcommand{\enorm}[1]{\lVert #1\rVert_{\mathrm{E}}}
\newcommand{\enormh}[1]{\lVert #1\rVert_{\mathrm{E},h}}
% Formulations: conventional spectral element, unsymmetric spectral element
% with the metric space Q_p in skew coordinates (USEM of the companion paper),
% unsymmetric spectral element with the hybrid trial space
\newcommand{\PP}{SEM}
\newcommand{\PM}{USEM}
\newcommand{\PH}{USEM-H}
\newcommand{\data}{data}   % path to the data files (relative to main.tex)

% Plot styles
\pgfplotscreateplotcyclelist{bases}{%
  {red,mark=*},{blue,mark=square*},{black,mark=triangle*},
  {green!50!black,mark=diamond*},{orange,mark=pentagon*}}
\pgfplotsset{
  every axis/.append style={tick label style={font=\footnotesize},
    label style={font=\small}, title style={font=\small},
    legend style={font=\footnotesize}},
  styPP/.style={black,mark=o,mark size=1.6pt,thick},
  styPM/.style={red,mark=square*,mark size=1.4pt,thick},
  styPH/.style={blue,mark=triangle*,mark size=1.8pt,thick},
  styLN/.style={black,dashed,thick},
}


% Styles of the thick-cylinder figures (colours for p = 2..8, dark = high p)
\newcommand{\eE}{e_{E}}
\definecolor{p2}{HTML}{FDB863}
\definecolor{p3}{HTML}{F08B3E}
\definecolor{p4}{HTML}{D9541A}
\definecolor{p5}{HTML}{B2182B}
\definecolor{p6}{HTML}{7F1C6A}
\definecolor{p7}{HTML}{4A2A8A}
\definecolor{p8}{HTML}{1B2A5E}
\definecolor{cPP}{HTML}{555555}
\definecolor{cPM}{HTML}{C0392B}
\definecolor{cPH}{HTML}{1F5FA8}
\pgfplotsset{
  conv/.style={
    xmode=log, ymode=log, x dir=reverse, grid=major,
    xtick={1,0.5,0.25,0.125}, xticklabels={$1$,$\frac12$,$\frac14$,$\frac18$},
    xmin=0.1, xmax=1.25, unbounded coords=jump, mark size=1.2pt, line width=0.7pt,
  },
  pline/.style={color=p#1, mark=*, mark options={solid,fill=p#1}},
  meshplot/.style={width=3.0cm, height=3.0cm, axis equal image, hide axis,
    title style={font=\scriptsize, yshift=-1.2ex}},
}

% Avoid overfull lines with long unbreakable math/macros
\setlength{\emergencystretch}{2em}

%
% Float placement: fuller float pages and more floats per text page, so that
% no float page is left largely empty
\renewcommand{\floatpagefraction}{0.85}
\setcounter{topnumber}{3}
\setcounter{totalnumber}{5}
%
% -----------------------------------------------------------------------------
% Title block
% -----------------------------------------------------------------------------
\title{\papertitle}
%\date{}                         % no date
%\date{September 9, 1985}        % fixed date
%
\author{
  \href{https://orcid.org/0000-0001-8774-9732}{\includegraphics[scale=0.06]{./logo/orcid.pdf}\hspace{1mm}S.~Eisentr\"ager\textsuperscript{1,$*$}}\\
  \href{mailto:sascha.eisentraeger@ovgu.de}{\texttt{sascha.eisentraeger@ovgu.de}}\\
  Faculty of Mechanical Engineering\\
  Institute of Mechanics\\
  Chair of Computational Mechanics\\
  Otto von Guericke University Magdeburg\\
  Magdeburg, 39106, Germany
%  % further authors to be added (side by side: \And; new row: \AND)
%  \And
%  \href{https://orcid.org/0000-0000-0000-0000}{\includegraphics[scale=0.06]{./logo/orcid.pdf}\hspace{1mm}F.~Surname\textsuperscript{2}}\\
%  \href{mailto:firstname.surname@ovgu.de}{\texttt{firstname.surname@ovgu.de}}\\
%  Faculty of Mechanical Engineering\\
%  Institute of Mechanics\\
%  Otto von Guericke University Magdeburg\\
%  Magdeburg, 39106, Germany
}
%
\renewcommand{\headeright}{Preprint submitted to \textnormal{\itshape\journalname}}
\renewcommand{\undertitle}{Preprint submitted to \textnormal{\itshape\journalname}}
\renewcommand{\shorttitle}{\papershorttitle}
%
% =============================================================================
\begin{document}
%
%\linenumbers   % line numbers, e.g. for the review version
%
\maketitle
%
\begin{abstract}
Unsymmetric (Petrov-Galerkin) spectral elements use conforming parametric test functions and metric trial functions that are complete in the physical coordinates. On meshes whose edge curvature persists under refinement, they therefore converge with the optimal rate, whereas the conventional spectral element of degree $p$ with a quadratic geometry is complete only to degree $\lfloor p/2\rfloor$. The metric trial functions are, however, obtained by nodal interpolation in the tensor-product polynomials of skew coordinates, and on strongly distorted elements they overshoot between the nodes, increasingly with the degree. In the paper at hand, we separate two causes of this behavior and propose one remedy for each. The first cause is the conditioning of the construction. Since the metric shape functions do not depend on the polynomial basis of the generalized Vandermonde matrix, we use the Lagrange polynomials through the reference Gauss-Lobatto-Legendre (GLL) points in skew coordinates scaled to the bounding box of the nodes. On parallelograms, the Vandermonde matrix is then a permutation matrix, and on distorted elements its condition number at $p=16$ is three to four orders of magnitude smaller than with monomials. Its condition number equals the Lebesgue function of the metric interpolation at the reference grid of the bounding box, so that no basis can reduce the Lebesgue constant itself. The second cause is the interpolation problem. Since exact reproduction requires only the complete polynomials of degree $p$ in the physical coordinates, we write only these in skew coordinates and the remaining functions, the $L^2$-orthogonal complement of the complete polynomials in the tensor-product space, in the reference coordinates, in which the nodes form the exact GLL grid. The resulting hybrid element coincides with the conventional spectral element on every straight-sided quadrilateral, and its unisolvence reduces to the regularity of a matrix of the size of the complete polynomial space. Its Lebesgue constant stays at or close to the GLL value for moderate edge curvature, where that of the tensor-product metric space reaches $10^4$--$10^7$ at $p=16$. On meshes with persistent edge curvature, the hybrid element converges with the optimal rate for $p\le6$, whereas the original unsymmetric element diverges for $p=5,6$. For a distorted thick cylinder, its error is up to 390 times smaller, and it is 0.78 to 1.00 times the interpolation error of its trial space on all meshes without folded elements. Moreover, the eigenvalues with negative real part of the original element are removed, although a slow growth of the transient response remains. The price is an error that is 1.3 to 5.5 times larger where the original element is stable and, on folded elements, singular integrands that limit the attainable accuracy.
\end{abstract}

%
\keywords{Unsymmetric finite element method \and Petrov-Galerkin method \and Spectral element method \and Curved-edge distortion \and Metric shape functions \and Lebesgue constant \and Generalized Vandermonde matrix}
%
\footnote[0]{$^*$Corresponding author. The superscript numbers indicate the order of the authors.}
\renewcommand*{\thefootnote}{\arabic{footnote}}
\setcounter{footnote}{0}
%
\section{Introduction}
\label{sec:intro}

Isoparametric finite elements lose accuracy on distorted meshes because the
isoparametric map destroys the completeness of the approximation space in
Cartesian coordinates. Under angular distortion, serendipity elements no longer
contain all quadratic polynomials \citep{Lee1993effects,Rajendran2000completeness},
and under curved-edge distortion even full tensor-product (Lagrange) elements of
degree $p$ with a geometry map of degree $g$ are guaranteed to contain $\Pp_k(x,y)$ only
for $gk\le p$ \citep[cf.][]{Arnold2002Approximation}.
MacNeal's theorem \citep{MacNeal1992On} shows that, within the class of symmetric
(Bubnov--Galerkin) elements, distortion insensitivity and passing the patch test
cannot be combined in general. To circumvent this restriction,
\citet{Rajendran2003novel} proposed the unsymmetric finite element method (UFEM), a
Petrov--Galerkin formulation in which test and trial functions differ. Here, the test
functions are the usual parametric (isoparametric) functions, which guarantee
inter-element compatibility. The trial functions, on the other hand, are \emph{metric}
functions, i.e.\ polynomials in Cartesian coordinates that are interpolatory at the
element nodes. \Citet{Rajendran2010technique} formulated three sufficient conditions for
\emph{distortion immunity}: (i) compatible test functions, (ii) trial functions that are
complete in Cartesian coordinates and (iii) an element stiffness integrand free of
$\det\matr J$ in the denominator. Elements satisfying them reproduce every displacement
field of the completeness degree exactly on arbitrarily distorted meshes, including
meshes with $\det\matr J\le0$. The approach has been extended to frame-invariant
formulations in skew coordinates \citep{Yuan1993New,Ooi2008Remedies,Xie2016Modified}, to
analytical trial functions \citep{Cen2012shape,Cen2015unsymmetric,Cen2019Some}, to
orthotropic and incompatible-mode elements \citep{Shang2018airy,Huang2022Unsymmetric}, to
solids \citep{Ooi2004hexahedron}, to enhanced assumed strain formulations
\citep{Pfefferkorn2021mesh,Pfefferkorn2023hourglassing}, to geometrically nonlinear
problems \citep{Ooi2007extension,Li2018high} and to elastodynamics
\citep{Liew2006quadratic,Wang2022Temporal,Zaehringer2026On}, to name just a few. The
non-variational character of the formulation was analysed by
\citet{Prathap2006mesh,Prathap2007The}.

Almost all of this work concerns low-order elements (4 to 20 nodes). In the article at
hand, we consider the generalisation to spectral elements of arbitrary degree $p$. Here,
the parametric test functions are Lagrange polynomials on Gauss--Lobatto--Legendre (GLL)
nodes, and the metric trial functions are the unique functions of a
$(p+1)^2$-dimensional space that interpolate the GLL nodes in physical space. In the
standard construction, the metric space is $\Qp_p(s,t)$, the tensor-product polynomials in
the non-dimensional corotational skew coordinates $(s,t)$ of \citet{Xie2016Modified}, and
the shape functions are obtained by inverting a generalised Vandermonde matrix built from
the monomials $s^at^b$ \citep{Rajendran2003novel,Ooi2006unsymmetric}. At high order,
however, two difficulties of a different nature appear:
\begin{enumerate}
\item \emph{Ill-conditioning of the construction.} The monomial Vandermonde matrix
  is exponentially ill-conditioned in $p$ already on undistorted elements
  \citep{Gautschi1988Lower}. Consequently, the computed shape functions carry round-off
  errors of order $\cond(\matr V)\,u$, where $u$ denotes the unit round-off. These errors
  destroy, for instance, the symmetry of the stiffness matrix on affine elements and the
  Kronecker property on which the diagonal GLL mass matrix relies.
\item \emph{Ill-conditioning of the interpolation problem itself.} On non-affine
  elements, the images of the GLL nodes do not form a tensor grid in $(s,t)$. The
  Lebesgue constant of the interpolation in $\Qp_p(s,t)$ then grows exponentially with
  $p$ and with the amount of distortion. Note that this is a property of the space and of
  the node set, independent of how the shape functions are computed. The Lebesgue
  constant amplifies round-off, enters the interpolation error and -- as we show in
  \cref{sec:numerics} -- is accompanied by a loss of stability of the Petrov--Galerkin
  operator. That is to say, strongly complex and negative generalised eigenvalues appear,
  which make the semi-discrete equations of motion exponentially unstable, and the static
  solution diverges under mesh refinement.
\end{enumerate}

Therefore, the overarching goal of this paper is to present two improvements of the
metric shape functions that address these two difficulties separately.

The first improvement concerns the polynomial basis of the metric Vandermonde matrix
(\cref{sec:basis}). We show that the metric shape functions are independent of the
polynomial basis used to construct them (\cref{prop:invariance}). Thus, the basis is a
purely numerical choice, and we propose the Lagrange polynomials through the reference
GLL nodes in bounding-box scaled skew coordinates. On affine elements, the Vandermonde
matrix then becomes a permutation matrix ($\cond=1$ for every $p$). On distorted
elements, its condition number equals the Lebesgue function of the metric interpolation
on the reference grid of the bounding box, up to a factor bounded by the reference
Lebesgue constant (\cref{prop:kappa}). This traces the remaining ill-conditioning back to
the extrapolation of the interpolant beyond the element. Numerically, the condition
number at $p=16$ drops by three to four orders of magnitude relative to monomials. The
same analysis also shows that \emph{no} choice of basis can reduce the Lebesgue
constant.

The second improvement is a hybrid metric--parametric trial space (\cref{sec:hybrid}).
Distortion immunity requires only that the trial space contain the complete polynomials
$\Pp_p(x,y)$, i.e.\ $(p+1)(p+2)/2$ of the $(p+1)^2$ functions. Therefore, we propose the
trial space
\[
  \Tp_p = \Pp_p(x,y)\big|_{\Omega_e}\;\oplus\;\Cp_p\circ\vect F_e^{-1},
\]
in which only the complete polynomials are written in metric (skew) coordinates. The
remaining approximation, $\Cp_p$, is the $L^2$-orthogonal complement of
$\Pp_p(\xi,\eta)$ in $\Qp_p(\xi,\eta)$, written in the parametric coordinates, in
which the nodes form an exact tensor grid. We prove that $\Tp_p$ is complete,
affine invariant and interpolatory and that it coincides with the classical spectral
element space on every straight-edged quadrilateral. Furthermore, we derive an exact
representation of its interpolation operator as the GLL interpolant plus a correction
that is driven only by the completeness defect of the classical element
(\cref{prop:representation}). A factorisation of its Vandermonde determinant reduces
unisolvence to the non-singularity of a $\dim\Pp_p\times\dim\Pp_p$ matrix
(\cref{prop:factor}). Numerically, its Lebesgue constant stays at the level of GLL
interpolation for moderate curved-edge distortion and remains orders of magnitude below
that of $\Qp_p(s,t)$ for strong distortion. On all meshes tested, the eigenvalues with
negative real part of the $\Qp_p(s,t)$ element disappear, and the relative imaginary
parts drop from $\mathcal O(1)$ to below $2\cdot10^{-2}$.

The numerical studies (\cref{sec:numerics}) quantify conditioning, accuracy of the
shape functions, Lebesgue constants versus degree and distortion, patch tests on
distorted and folded meshes, $h$-convergence on meshes whose distortion persists
under refinement, eigenvalue spectra, transient stability and the thick-cylinder
benchmark of \citet{Rajendran2010technique}. In order to make the trade-off explicit,
they also document where the hybrid space is \emph{less} accurate than $\Qp_p(s,t)$
(static problems with smooth solutions at moderate degree and distortion). Note that only
linear plane elasticity on quadrilaterals is considered.

The outline of the remainder of the present article is as follows: First, the
unsymmetric spectral element formulation and the metric shape functions in skew
coordinates are briefly summarised in \cref{sec:formulation}. Thereafter, the polynomial
basis of the metric Vandermonde matrix is discussed in \cref{sec:basis}, before the
hybrid trial space is introduced and analysed in \cref{sec:hybrid}. The implementation is
outlined in \cref{sec:implementation}, and the numerical studies are presented in
\cref{sec:numerics}. Finally, the results are discussed in \cref{sec:discussion}, and the
most important conclusions are summarised in \cref{sec:conclusions}.

\section{Unsymmetric spectral elements}
\label{sec:formulation}

\subsection{Problem and Petrov--Galerkin discretisation}

We consider linear plane elasticity on $\Omega\subset\R^2$,
\begin{equation}
  -\nabla\cdot\bm\sigma(\vect u) = \vect b \ \text{in } \Omega,\qquad
  \bm\sigma = \matr C\,\bm\varepsilon(\vect u),\qquad
  \vect u = \bar{\vect u}\ \text{on } \Gamma_u,\qquad
  \bm\sigma\vect n = \bar{\vect t}\ \text{on } \Gamma_t ,
  \label{eq:strong}
\end{equation}
with the strain vector $\bm\varepsilon=[\varepsilon_{xx},\varepsilon_{yy},
\gamma_{xy}]^\T$ and the plane stress or plane strain matrix $\matr C$. The mesh
consists of quadrilaterals $\Omega_e=\vect F_e(\hat\Omega)$,
$\hat\Omega=[-1,1]^2$, with an isoparametric geometry map $\vect F_e$ of degree $p$
in each reference coordinate $\bxi=(\xi,\eta)$. The element nodes are the images
$\bx_k=\vect F_e(\bxi_k)$, $k=1,\dots,n=(p+1)^2$, of the tensor-product GLL nodes
$\bxi_k$; for simplicity of notation $p_\xi=p_\eta=p$ unless stated otherwise.

The \emph{parametric} functions $N_k=\hat N_k\circ\vect F_e^{-1}$, where $\hat
N_k\in\Qp_p(\xi,\eta)$ are the Lagrange polynomials through the GLL nodes, span the
conforming test space $\mathcal W_h\subset H^1(\Omega)$. The \emph{metric} trial
functions $M_k^e$ span an element space $\mathcal U_e$ of dimension $n$ and are
interpolatory, $M_k^e(\bx_l)=\delta_{kl}$; the global trial function is
$\vect u_h|_{\Omega_e}=\sum_k \hat{\vect u}_k M_k^e$ with the nodal displacements
$\hat{\vect u}_k$ shared by adjacent elements. The trial functions are in general
discontinuous across element edges (they coincide only at the nodes). The
discrete problem reads: find the nodal values such that
\begin{equation}
  \sum_e\int_{\Omega_e}\bm\varepsilon(\vect w_h)^\T\matr C\,\bm\varepsilon(\vect u_h)\dd\Omega
  = \int_\Omega\vect w_h^\T\vect b\dd\Omega + \int_{\Gamma_t}\vect w_h^\T\bar{\vect t}\dd\Gamma
  \quad\forall\,\vect w_h\in\mathcal W_h,\ \vect w_h|_{\Gamma_u}=\vect 0 .
  \label{eq:pg}
\end{equation}
In matrix form the element stiffness is
\begin{equation}
  \matr K_e=\int_{\hat\Omega}\tilde{\matr B}_L^\T\,\matr C\,\matr B_R\dd\xi\dd\eta,
  \qquad \tilde{\matr B}_L=\det\matr J\,\matr B_L ,
  \label{eq:Ke}
\end{equation}
where $\matr B_L$ contains the Cartesian derivatives of the parametric functions and
$\matr B_R$ those of the metric functions, evaluated at $\vect F_e(\bxi)$.
$\tilde{\matr B}_L$ is assembled from $\adj\matr J\,\nabla_{\bxi}\hat N_k$, so
$\det\matr J$ does not appear in a denominator. $\matr K_e$ is unsymmetric unless
$\mathcal U_e=\mathcal W_h|_{\Omega_e}$. The consistent mass matrix is
$\matr M_e=\int_{\hat\Omega}\rho\,\matr N^\T\matr M\det\matr J\dd\xi\dd\eta$ with
the parametric ($\matr N$) and metric ($\matr M$) function values; with GLL nodal
quadrature it is diagonal, because both families are interpolatory at the GLL nodes.
We refer to the formulation as \PP\ (parametric test and trial functions, i.e.\ the
classical spectral element) or \PM\ (parametric test, metric trial functions), and to
\PM\ with the hybrid trial space of \cref{sec:hybrid} as \PH.

\subsection{Distortion immunity}

\begin{proposition}[Distortion immunity]
\label{prop:immunity}
Let every element map $\vect F_e$ be bijective, let $\Pp_q(x,y)\subset\mathcal U_e$ for
every element, let the exact solution of \eqref{eq:strong} satisfy
$\vect u\in[\Pp_q(x,y)]^2$, let the quadrature of \eqref{eq:Ke} and of the load vector be
exact for the resulting integrands, and let the global system matrix (after imposing the
Dirichlet conditions) be non-singular. Then the solution of \eqref{eq:pg} is
$\vect u_h=\vect u$, independently of the shape of the elements.
\end{proposition}
\begin{proof}
Because $\mathcal U_e$ is interpolatory and contains $\vect u|_{\Omega_e}$, the
interpolant of the exact nodal values is $\vect u$ itself on every element, so
$\bm\varepsilon(\vect u_h)=\bm\varepsilon(\vect u)$. As $\vect w_h$ is continuous,
integration by parts over $\Omega$ gives
$\int_\Omega\bm\varepsilon(\vect w_h)^\T\bm\sigma(\vect u)\dd\Omega
=\int_\Omega\vect w_h^\T\vect b\dd\Omega+\int_{\Gamma_t}\vect w_h^\T\bar{\vect t}\dd\Gamma$.
Hence the exact nodal values satisfy \eqref{eq:pg}, and they are its unique solution.
\end{proof}

\begin{remark}[Folded elements]
\label{rem:folded}
If $\vect F_e$ is not injective (``folded'' elements, $\det\matr J$ changing sign), the
parametric functions are not well defined in physical space and the argument above does
not apply. The element matrices \eqref{eq:Ke}, however, are defined in reference
coordinates with the signed $\det\matr J$ (through $\adj\matr J$), as are the loads. For
such elements exactness of polynomial solutions has been observed for low-order elements
\citep{Rajendran2010technique} and is confirmed by our patch tests
(\cref{sec:num-patch}); we do not give a proof here.
\end{remark}
This is the higher-order form of the conditions of \citet{Rajendran2010technique}.
The parametric space violates the hypothesis on curved elements: for a geometry map
of degree $g$ in each variable, $x^ay^b\circ\vect F_e\in\Qp_{g(a+b)}$, so that
$\Pp_k(x,y)\circ\vect F_e\subset\Qp_p(\xi,\eta)$ is guaranteed only for $gk\le p$
\citep[cf.][]{Arnold2002Approximation}. For bilinear maps ($g=1$), the classical
tensor-product element is complete of degree $p$; the loss of completeness under
purely angular distortion reported by \citet{Lee1993effects} concerns serendipity
spaces.

\subsection{Metric shape functions in skew coordinates}
\label{sec:skew}

Following \citet{Yuan1993New} and \citet{Xie2016Modified}, the metric functions are
written in the corotational skew coordinates
\begin{equation}
  \vect s=(s,t)^\T=\matr J_0^{-\T}(\bx-\bx_0),\qquad
  \bx_0=\vect F_e(\vect 0),\quad
  \matr J_0=\frac{\partial(x,y)}{\partial(\xi,\eta)}\Big|_{\bxi=\vect 0},
  \label{eq:skew}
\end{equation}
which require $\det\matr J_0\ne0$, with the Jacobian convention
$\matr J=\bigl[\begin{smallmatrix}x_{,\xi}&y_{,\xi}\\x_{,\eta}&y_{,\eta}\end{smallmatrix}\bigr]$.

\begin{lemma}[Affine invariance]
\label{lem:affine}
Let $\bx\mapsto\matr A\bx+\vect c$ with $\det\matr A\ne0$ be applied to all nodes of
an element. Then the skew coordinates of every material point are unchanged. In
particular, any space defined in $(s,t)$ (or in $(\xi,\eta)$) is invariant under
rigid-body motions, uniform scaling and every other affine transformation of the
element.
\end{lemma}
\begin{proof}
The transformed map is $\vect F_e'=\matr A\vect F_e+\vect c$, hence
$\matr J_0'^\T=\matr A\matr J_0^\T$ and $\bx'-\bx_0'=\matr A(\bx-\bx_0)$, so
$\vect s'=\matr J_0^{-\T}\matr A^{-1}\matr A(\bx-\bx_0)=\vect s$.
\end{proof}
Because $\vect s$ is affine in $\bx$, $\Qp_p(s,t)\supset\Pp_{p}(s,t)=\Pp_p(x,y)$, so
the space $\mathcal U_e=\Qp_p(s,t)$ satisfies the hypothesis of \cref{prop:immunity}
with $q=p$ on every element. On an affine element (parallelogram), $\vect F_e$ is
affine and $\vect s=\bxi$, so the metric and parametric functions coincide. Tensor
products of monomials in centroid-shifted Cartesian coordinates
\citep{Rajendran2003novel} span $\Qp_p(x,y)$, which contains $\Pp_p(x,y)$ as well but
depends on the orientation of the Cartesian frame (e.g.\ $xy$ versus $st$); this
makes the element frame dependent \citep{Ooi2008Remedies}.

Given a basis $\{\phi_c\}_{c=1}^n$ of $\mathcal U_e$, the metric functions are
computed from the generalised Vandermonde matrix and the basis values at the query
points $\bx_q$ (e.g.\ the images of the quadrature points),
\begin{equation}
  V_{ck}=\phi_c(\bx_k),\qquad P_{cq}=\phi_c(\bx_q),\qquad
  \matr N=\matr V^{-1}\matr P,\qquad N_{lq}=M_l^e(\bx_q),
  \label{eq:vandermonde}
\end{equation}
and analogously for the derivatives. Throughout, $\Lam_N$ denotes the Lebesgue
constant of tensor-product GLL interpolation on $\hat\Omega$, and for an
interpolatory space $\mathcal U_e$ with shape functions $M_l^e$
\begin{equation}
  \lambda_{\mathcal U}(\bx)=\sum_{l=1}^n|M_l^e(\bx)|,\qquad
  \Lam_{\mathcal U}=\max_{\bx\in\Omega_e}\lambda_{\mathcal U}(\bx)
  \label{eq:lebesgue}
\end{equation}
are its Lebesgue function and Lebesgue constant. $\Lam_{\mathcal U}=\|\mathcal
I_{\mathcal U}\|_{L^\infty\to L^\infty}$ is the norm of the nodal interpolation
operator $\mathcal I_{\mathcal U}$, so that
$\|f-\mathcal I_{\mathcal U}f\|_\infty\le(1+\Lam_{\mathcal U})\inf_{v\in\mathcal
U_e}\|f-v\|_\infty$. Note that the parametric Lebesgue constant is independent of the
element map, because $\sum_k|N_k(\vect F_e(\bxi))|=\sum_k|\hat N_k(\bxi)|$.

\section{Improvement 1: the basis of the metric Vandermonde matrix}
\label{sec:basis}
%
\subsection{Basis invariance and floating-point error}
\label{sec:basis-invariance}
%
The point of departure is the observation that the metric shape functions are defined by the space and the nodes alone.

\begin{proposition}[Basis invariance]
\label{prop:invariance}
If $\matr V$ in \cref{eq:vandermonde} is regular, the metric functions are the unique functions $M_i^K\in\mathcal U_K$ with $M_i^K(\bx_j)=\delta_{ij}$. They do not depend on the basis $\{\phi_c\}$.
\end{proposition}
\begin{proof}
For another basis $\bm\psi=\matr G\bm\phi$ with a regular matrix $\matr G$, we obtain $\matr V_\psi=\matr G\matr V_\phi$ and $\matr P_\psi=\matr G\matr P_\phi$, so that $\matr V_\psi^{-1}\matr P_\psi=\matr V_\phi^{-1}\matr P_\phi$. Uniqueness follows from the regularity of $\matr V$.
\end{proof}

Completeness, affine invariance, exact reproduction and unisolvence are therefore properties of the space and the node set, and so is the Lebesgue constant \cref{eq:lebesgue}. That is to say, the basis affects only the floating-point error. For the solution of $\matr V\matr N_{\mathrm M}=\matr P$ by Gaussian elimination with partial pivoting, the standard perturbation bound \citep{Higham2002Accuracy} gives
%
\begin{equation}
  \frac{\|\hat{\matr N}_{\mathrm M}-\matr N_{\mathrm M}\|}{\|\matr N_{\mathrm M}\|}\lesssim c_n\,\rho_n\,\cond(\matr V)\,u ,
  \label{eq:fp}
\end{equation}
%
where $\hat{\matr N}_{\mathrm M}$ is the computed matrix, $\rho_n$ the growth factor, which is moderate in practice, $c_n$ a constant of moderate size and $u\approx1.1\cdot10^{-16}$ the unit round-off. Thus, the basis has to be chosen such that $\cond(\matr V)$ is small.
%
\subsection{Candidate bases and bounding-box scaling}
\label{sec:basis-candidates}
%
To this end, we consider tensor-product bases $\phi_c=\varphi_a(\tilde\xi)\,\varphi_b(\tilde\eta)$, $0\le a,b\le p$, with four choices of the one-dimensional factors: (i) the monomials $\varphi_a(\tilde\xi)=\tilde\xi^a$ of the original construction, (ii) the Legendre polynomials $\varphi_a=L_a$, (iii) the Chebyshev polynomials $\varphi_a=T_a$ and (iv) the Lagrange polynomials $\varphi_a=\ell_a$ through the $p+1$ GLL points of $[-1,1]$, i.e., the one-dimensional factors of the reference shape functions. The Legendre and Chebyshev polynomials are evaluated by their three-term recurrences and the Lagrange polynomials by the product formula. The last three bases are designed for the interval $[-1,1]$. On distorted elements, however, the skew coordinates of the nodes leave $[-1,1]^2$, and outside this interval the polynomials grow rapidly, e.g., $|T_p(\bar\xi)|=\cosh\bigl(p\operatorname{arcosh}|\bar\xi|\bigr)$ for $|\bar\xi|>1$. Therefore, we map the bounding box of the nodes in the skew coordinates onto $[-1,1]^2$,
%
\begin{equation}
  \begin{aligned}
  &\tilde\xi=\frac{\bar\xi-c_{\xi}}{r_{\xi}}\qquad\text{with}\\
  &c_{\xi}=\tfrac12\bigl(\max_j\bar\xi_j+\min_j\bar\xi_j\bigr)\quad\text{and}\quad
  r_{\xi}=\tfrac12\bigl(\max_j\bar\xi_j-\min_j\bar\xi_j\bigr),
  \end{aligned}
  \label{eq:bbox}
\end{equation}
%
and likewise for $\tilde\eta$, where $\bar\xi_j$ is the first skew coordinate of node $j$. The map is affine and computed from affine-invariant quantities, so that the space $\Qsk$, its completeness and the invariance of \cref{lem:affine} are preserved. On a parallelogram, the skew coordinates of the nodes are the GLL points $\bxi_j$, whose bounding box is $\Kh$, and \cref{eq:bbox} reduces to the identity. The derivatives follow from $\nabla_{\bx}=\matr J_0^{-1}\operatorname{diag}(r_{\xi}^{-1},r_{\eta}^{-1})\,\nabla_{\bxit}$, where $\bxit=(\tilde\xi,\tilde\eta)$ denotes the scaled skew coordinates.
%
\subsection{Parallelograms}
\label{sec:basis-affine}
%
On a parallelogram, the nodes form the GLL tensor grid in $(\tilde\xi,\tilde\eta)$. Thus, $\matr V$ is, up to a permutation of the columns, the Kronecker product $\matr V_1\otimes\matr V_1$ of the one-dimensional Vandermonde matrix, and $\cond_2(\matr V)=\cond_2(\matr V_1)^2$. For monomials on nodes that are symmetric about the origin, $\cond(\matr V_1)$ grows at least exponentially \citep{Gautschi1988Lower}, so that the condition number of $\matr V$ grows with the square of an exponential. For the Lagrange basis, on the other hand, $\phi_c(\bxit_j)=\delta_{cj}$ holds up to the ordering, which leads to the following result.

\begin{corollary}
\label{cor:affine}
On a parallelogram, the Lagrange basis gives a permutation matrix $\matr V$ (up to the rounding of the scaled skew coordinates of the nodes), $\cond_r(\matr V)=1$ in every $r$-norm, and $\matr N_{\mathrm M}=\matr V^{-1}\matr P$ is obtained without amplification of rounding errors. The metric and the parametric functions then agree to round-off, and $\matr K_K$ is symmetric to round-off, for every $p$.
\end{corollary}
%
\subsection{Distorted elements}
\label{sec:basis-distorted}
%
On a non-affine element, the condition number of $\matr V$ in the Lagrange basis has a precise interpretation. Let $\bm g_d$, $d=1,\dots,n$, denote the GLL tensor grid placed in the plane of the scaled skew coordinates, so that $\phi_c(\bm g_d)=\delta_{cd}$, and let $\lambda_{\mathrm{ref}}(\bxit)=\sum_c|\phi_c(\bxit)|$ be the Lebesgue function of the GLL interpolation in these coordinates.

\begin{proposition}[Condition number in the Lagrange basis]
\label{prop:kappa}
With the Lagrange basis and the scaled skew coordinates $\bxit_j$ of the element nodes,
%
\begin{equation}
  (\matr V^{-1})_{id}=M_i^K(\bm g_d)\qquad\text{and}\qquad
  \cond_1(\matr V)=\max_j\lambda_{\mathrm{ref}}(\bxit_j)\;\max_d\lambda_{\mathcal U}(\bm g_d) ,
  \label{eq:kappa}
\end{equation}
%
where $M_i^K$ and $\lambda_{\mathcal U}$ are understood as functions of the scaled skew coordinates.
\end{proposition}
\begin{proof}
Choosing the query points $\bm g_d$ in \cref{eq:vandermonde} gives $\matr P=\matr I$ and hence $\matr N_{\mathrm M}=\matr V^{-1}$, which is the first identity. The 1-norm is the largest absolute column sum, i.e., $\|\matr V\|_1=\max_j\sum_c|\phi_c(\bxit_j)|$ and $\|\matr V^{-1}\|_1=\max_d\sum_i|M_i^K(\bm g_d)|$.
\end{proof}

By \cref{eq:bbox}, all $\bxit_j$ lie in $\Kh$, and the first factor in \cref{eq:kappa} is bounded by $\Lam_N$ ($\Lam_N\le6.1$ for $p\le16$). The second factor is the Lebesgue function of the metric interpolation at the reference grid of the bounding box. On a parallelogram, this grid coincides with the nodes, and the factor is one. On a distorted element, however, part of the grid lies \emph{outside} the element, e.g., the corners of the bounding box. There, $\lambda_{\mathcal U}$ is an extrapolation and can exceed the Lebesgue constant $\Lam_{\mathcal U}$ by orders of magnitude (\cref{sec:num-conditioning}). Thus, \cref{prop:kappa} traces the remaining ill-conditioning of every tensor-product basis on the bounding box back to the extrapolation of the metric interpolant from the element to its bounding box. This effect cannot be removed by switching among the bases considered here, which numerically lie within a factor of about 3 of one another. It can only be removed by bases adapted to the node set itself, e.g., discretely orthogonalized bases \citep{Brubeck2021Vandermonde}, which is out of the scope of this contribution.

The proposition also shows the limit of \emph{any} basis. By \cref{prop:invariance}, $\lambda_{\mathcal U}$ and $\Lam_{\mathcal U}$ do not depend on the basis. Therefore, a large Lebesgue constant amplifies perturbations of the nodal values, i.e., round-off or the nodal error of a discrete solution, by up to $\Lam_{\mathcal U}$, whatever basis is used. Reducing $\Lam_{\mathcal U}$ requires a different \emph{space}, which is the subject of \cref{sec:hybrid}.

Note that the effect of $\cond(\matr V)$ depends on the nodal data. For the nodal data $\vect f=\matr V^\T\vect c$ of a function in $\mathcal U_K$ with the coefficient vector $\vect c$, the first-order error of the computed interpolant is $\vect c^\T\delta\matr V\,\matr N_{\mathrm M}$, where $\delta\matr V$ denotes the backward error of the factorization. Up to a constant, this error is bounded by $u\,\|\vect c\|\,\|\matr V\|\,\Lam_{\mathcal U}$. For polynomials of low degree, $\|\vect c\|\,\|\matr V\|$ is moderate in every basis. Consequently, the reproduction of polynomials, and with it the patch test, is insensitive to $\cond(\matr V)$ and limited by $u\,\Lam_{\mathcal U}$. Generic nodal data such as unit vectors, eigenvectors or wave fields, on the other hand, can have $\|\vect c\|\approx\|\matr V^{-1}\|\,\|\vect f\|$, and the error is then of the order of $\cond(\matr V)\,u$. Since the element matrices are built from the individual functions $M_i^K$, i.e., from unit nodal data, they inherit the bound \cref{eq:fp}.

\section{Improvement 2: a hybrid metric-parametric trial space}
\label{sec:hybrid}
%
\subsection{Requirements and freedom in the choice of the space}
\label{sec:hybrid-requirements}
%
A metric trial space $\mathcal U_K$ for the formulation of \cref{sec:formulation} has to satisfy three requirements:
\begin{enumerate}
\item[(R1)] $\dim\mathcal U_K=n=(p_\xi+1)(p_\eta+1)$ and unisolvence at the element nodes, so that the shape functions are interpolatory, the nodal unknowns are shared and the GLL mass matrix is diagonal;
\item[(R2)] $\Pp_m(x,y)\subset\mathcal U_K$ with $m=\min(p_\xi,p_\eta)$, i.e., completeness in the sense of condition (C1) of \cref{sec:exact}, which includes the rigid-body modes;
\item[(R3)] invariance under rigid-body motions and scaling of the element (frame invariance).
\end{enumerate}
Requirement (R2) fixes $(m+1)(m+2)/2$ of the $n$ dimensions. For $p_\xi=p_\eta=p$, the remaining $p(p+1)/2$ dimensions are free. The space $\Qsk$ fills them with Cartesian polynomials of total degree $p+1,\dots,2p$. On a non-affine element, the nodes do not form a tensor-product grid in $(\bar\xi,\bar\eta)$, and we attribute the growth of the Lebesgue constant (\cref{sec:num-lebesgue}) to the interpolation of these Cartesian polynomials of high degree on the perturbed grid (cf.\ \cref{prop:representation} below). In the reference coordinates, on the other hand, the nodes form the exact GLL tensor grid on every element. This motivates our choice to write only the complete polynomials in the skew coordinates and the remaining functions in the reference coordinates.
%
\subsection{Definition}
\label{sec:hybrid-definition}
%
\begin{definition}[Parametric complement and hybrid space]
Let $L_a$ denote the Legendre polynomials and
%
\begin{equation}
  \Cp_p=\spn\bigl\{L_a(\xi)L_b(\eta):\ 0\le a\le p_\xi,\ 0\le b\le p_\eta,\ a+b>m\bigr\}.
  \label{eq:complement}
\end{equation}
%
The hybrid trial space of element $K$ is defined through its pull-back to $\Kh$,
%
\begin{equation}
  \hat\Tp_p=\bigl\{q\circ\bx_K:\ q\in\Pp_m(x,y)\bigr\}+\Cp_p\qquad\text{and}\qquad
  \Tp_p=\bigl\{\hat v\circ\bx_K^{-1}:\ \hat v\in\hat\Tp_p\bigr\},
  \label{eq:hybrid}
\end{equation}
%
and its shape functions are the functions $\hat M_i\in\hat\Tp_p$ with $\hat M_i(\bxi_j)=\delta_{ij}$.
\end{definition}

The complete polynomials are represented by $L_a(\tilde\xi)L_b(\tilde\eta)$, $a+b\le m$, in the scaled skew coordinates \cref{eq:bbox}. The definition through $\Kh$ requires only the evaluation of $\bx_K$, not of its inverse. Where $\bx_K$ is not injective (folded elements), only the pull-back $\hat\Tp_p$ is defined, whereas the physical space $\Tp_p$ and point evaluation at physical coordinates are not, and all element quantities are computed in the reference coordinates as in \cref{eq:Ke}.

\begin{lemma}[Canonical complement]
\label{lem:complement}
$\Cp_p$ is the $L^2(\Kh)$-orthogonal complement of $\Pp_m(\xi,\eta)$ in $\Qp_{p_\xi,p_\eta}(\xi,\eta)$. In particular, $\hat\Tp_p$ is defined independently of any basis.
\end{lemma}
\begin{proof}
Let $a+b>m$ and $i+j\le m$. Then $i<a$ or $j<b$, and the orthogonality of $L_a$ to all polynomials of degree less than $a$ gives $\int_{\Kh}L_a(\xi)L_b(\eta)\,\xi^i\eta^j\dd\xi\dd\eta=\int L_a\xi^i\dd\xi\int L_b\eta^j\dd\eta=0$. The functions $L_aL_b$ are linearly independent, and $\dim\Cp_p=n-\dim\Pp_m$.
\end{proof}

For a non-affine map, $\Pp_m(x,y)\circ\bx_K\ne\Pp_m(\xi,\eta)$, so that a different complement of $\Pp_m(\xi,\eta)$, e.g., $\spn\{\xi^a\eta^b:a+b>m\}$, would give a different space $\hat\Tp_p$. \Cref{lem:complement} identifies \cref{eq:complement} as the canonical choice. Note that, unlike for $\Qsk$, the complement is therefore part of the definition of the space.
%
\subsection{Properties}
\label{sec:hybrid-properties}
%
\begin{proposition}[Properties of the hybrid space]
\label{prop:hybrid}
Let the Vandermonde matrix of the hybrid space be regular. Then
\begin{enumerate}
\item[(i)] $\dim\Tp_p=n$, the sum in \cref{eq:hybrid} is direct, and the shape functions exist and are unique;
\item[(ii)] $\Pp_m(x,y)\subset\Tp_p$, so that condition (C1) holds for all polynomial solutions of degree $m$, including the rigid-body modes;
\item[(iii)] $\Tp_p$ is invariant under affine transformations of the element (\cref{lem:affine});
\item[(iv)] if $\bx_K$ is bilinear (straight-sided quadrilateral), then $\Tp_p=\Qp_{p_\xi,p_\eta}(\xi,\eta)\circ\bx_K^{-1}$, i.e., the hybrid shape functions coincide with the parametric ones, and \PH\ is the conventional, symmetric spectral element.
\end{enumerate}
\end{proposition}
\begin{proof}
(i) $\hat\Tp_p$ is spanned by $n$ functions whose values at the $n$ nodes form the regular matrix $\matr V$. Hence, they are linearly independent, and the nodal map is bijective. (ii) and (iii) follow immediately from \cref{eq:hybrid}, \cref{lem:affine} and the invariance of $\bxi$. (iv) For a bilinear map, $x,y\in\Qp_1(\xi,\eta)$, and hence $x^ay^b\circ\bx_K\in\Qp_{a+b}(\xi,\eta)\subset\Qp_{p_\xi,p_\eta}(\xi,\eta)$ for $a+b\le m$. With $\Cp_p\subset\Qp_{p_\xi,p_\eta}(\xi,\eta)$, we obtain $\hat\Tp_p\subset\Qp_{p_\xi,p_\eta}(\xi,\eta)$, and equality follows from (i).
\end{proof}

Property (iv) means that the unsymmetric formulation becomes active only where it is needed, namely on elements with curved edges. On straight-sided meshes, \PH\ inherits the symmetry, the positive definiteness and the spectral properties of the conventional element, however strong the angular distortion may be. Since the conventional spectral element keeps its completeness under angular distortion (see \cref{sec:completeness}), nothing is lost on such meshes. The space $\Qsk$ does not have this property. In the companion paper, it was exactly on straight-sided elements with interior angles close to $180^\circ$ that the metric functions of \PM\ overshot most \cite{Eisentraeger2026axisymmetric}.
%
\subsection{Representation of the interpolation operator and Lebesgue constant}
\label{sec:hybrid-representation}
%
Let $\Ihat$ denote the GLL interpolation on $\Kh$ onto $\Qp=\Qp_{p_\xi,p_\eta}(\xi,\eta)$, $\pi$ the projector of $\Qp$ onto $\Pp_m(\xi,\eta)$ along $\Cp_p$ (truncation of the Legendre expansion), $\hat\Pp_m=\{q\circ\bx_K:q\in\Pp_m(x,y)\}$ and $\hat\Qp_p=\{q\circ\bx_K:q\in\Qsk\}$. We define
%
\begin{equation}
  S=\pi\Ihat\big|_{\hat\Pp_m}:\ \hat\Pp_m\to\Pp_m(\xi,\eta)\qquad\text{and}\qquad
  S_Q=\Ihat\big|_{\hat\Qp_p}:\ \hat\Qp_p\to\Qp ,
  \label{eq:S}
\end{equation}
%
and denote by $\mathcal I_T$ and $\mathcal I_Q$ the nodal interpolation operators of $\hat\Tp_p$ and $\hat\Qp_p$, all functions being pulled back to $\Kh$.

\begin{proposition}[Representation of the metric interpolants]
\label{prop:representation}
Assume that $\bx_K$ has a point with $\det(\matr J)\ne0$, so that $\dim\hat\Pp_m=\dim\Pp_m$.
\begin{enumerate}
\item[(a)] $\hat\Tp_p$ is unisolvent at the GLL points if and only if $S$ is invertible, and then, for every $f\in C(\Kh)$,
  \begin{equation}
    \mathcal I_Tf=\Ihat f+(I-\Ihat)\,S^{-1}\pi\,\Ihat f .
    \label{eq:repT}
  \end{equation}
\item[(b)] $\hat\Qp_p$ is unisolvent if and only if $S_Q$ is invertible, and then
  \begin{equation}
    \mathcal I_Qf=\Ihat f+(I-\Ihat)\,S_Q^{-1}\,\Ihat f .
    \label{eq:repQ}
  \end{equation}
\item[(c)] Consequently, with operator norms on $(\Qp,\|\cdot\|_{L^\infty(\Kh)})$,
  \begin{equation}
    \Lam_T\le\Lam_N\bigl(1+\|(I-\Ihat)S^{-1}\pi\|\bigr)\qquad\text{and}\qquad
    \Lam_Q\le\Lam_N\bigl(1+\|(I-\Ihat)S_Q^{-1}\|\bigr),
    \label{eq:lebbound}
  \end{equation}
  and the shape functions are $\hat M_{i,T}=\hat N_i+(I-\Ihat)S^{-1}\pi\hat N_i$ and $\hat M_{i,Q}=\hat N_i+(I-\Ihat)S_Q^{-1}\hat N_i$.
\end{enumerate}
\end{proposition}
\begin{proof}
The nodal values of a function $\hat v$ are in one-to-one correspondence with $\Ihat\hat v\in\Qp$, and $\Ihat c=c$ for $c\in\Cp_p\subset\Qp$. (a) Let $\hat v=q+c$ with $q\in\hat\Pp_m$ and $c\in\Cp_p$ have the nodal values of $f$, i.e., $\Ihat q+c=\Ihat f$. Applying $\pi$, whose kernel is $\Cp_p$, gives $Sq=\pi\Ihat f$. Conversely, for every $q$ with $Sq=\pi\Ihat f$, the function $c=\Ihat f-\Ihat q$ lies in $\ker\pi=\Cp_p$. If $S$ is invertible, $q$ and $c$ are therefore uniquely determined for every $f$, and $\hat v=q+\Ihat f-\Ihat q$, which is \cref{eq:repT}. If $S$ is singular, we take $0\ne q\in\ker S$ and $c=-\Ihat q\in\Cp_p$. The function $\hat v=q+c$ then vanishes at all nodes. Either $\hat v\ne0$, or $q=-c\in\hat\Pp_m\cap\Cp_p$, so that the sum is not direct and $\dim\hat\Tp_p<n$. In both cases, $\hat\Tp_p$ is not unisolvent. (b) An interpolant $\hat v\in\hat\Qp_p$ satisfies $\Ihat\hat v=\Ihat f$, i.e., $S_Q\hat v=\Ihat f$, and $\hat v=\Ihat\hat v+(I-\Ihat)\hat v$. (c) follows from $\|\Ihat f\|\le\Lam_N\|f\|$ and, for the shape functions, from $\Ihat f=\hat N_i$ for the nodal data $\delta_{\cdot i}$.
\end{proof}

\begin{corollary}
\label{cor:defect}
If $\hat\Pp_m\subset\Qp$, in particular for every bilinear map, the correction in \cref{eq:repT} vanishes, $\mathcal I_T=\Ihat$ and $\Lam_T=\Lam_N$, independently of the amount of angular distortion. In general, the deviation of $\mathcal I_T$ from the GLL interpolation is the GLL interpolation error of a complete polynomial of degree $m$, $(I-\Ihat)(q\circ\bx_K)$ with $q\in\Pp_m(x,y)$, i.e., precisely the \emph{completeness defect} of the conventional element.
\end{corollary}

\Cref{prop:representation} puts both metric spaces into the same framework, namely the GLL interpolation plus a correction, and it isolates the structural difference between them. For $\Tp_p$, the correction has at most the rank $\dim\Pp_m=(m+1)(m+2)/2$ and involves only the inverse $S^{-1}$ of the low-order block, so that the high-order part of the interpolant remains untouched. For $\Qsk$, in contrast, the correction acts on the whole $n$-dimensional space. It involves the GLL interpolation error of Cartesian polynomials of degree $p$ in each skew coordinate, i.e., of total degree up to $2p$, as well as the inverse of the full matrix $S_Q$. In particular, this correction does not vanish on bilinear elements, since in general $\bar\xi^a\bar\eta^b\circ\bx_K\in\Qp_{a+b}(\xi,\eta)\not\subset\Qp_p(\xi,\eta)$ for $a+b>p$. For a smooth perturbation $\bx_K=\bx_{\mathrm{bil}}+\delta\,\tilde{\bx}$ of a bilinear map, on the other hand, the correction of $\Tp_p$ is of the order of $\delta$, since $q\circ\bx_{\mathrm{bil}}\in\Qp_p(\xi,\eta)$ for $q\in\Pp_m(x,y)$. Note that the bounds \cref{eq:lebbound} are not sharp enough to quantify the difference a priori. The numerical results of \cref{sec:num-lebesgue} show that $\Lam_T$ stays at or very close to $\Lam_N$ for moderate edge curvature. For strong distortion, $\Lam_T$ grows as well, but by orders of magnitude less than $\Lam_Q$, which grows exponentially in $p$ and in the distortion.

The operator $S$ has a simple matrix form. Let $\matr V_{\mathrm{ref}}$ be the Legendre Vandermonde matrix of the GLL grid with the rows $L_a(\xi_j)L_b(\eta_j)$, $a\le p_\xi$, $b\le p_\eta$. Its rows are ordered into the $n_m=\dim\Pp_m$ low-order rows $\matr L$ ($a+b\le m$) and the rows $\matr R$ of the complement. The hybrid Vandermonde matrix is ordered alike, $\matr V_T=[\matr Q;\matr R]$, where the rows of $\matr Q$ are the nodal values of a basis $\{q_k\}$ of $\hat\Pp_m$. Let $\matr S\in\R^{n_m\times n_m}$ be the matrix of $S$, i.e., $S_{lk}$ is the coefficient of the $l$-th low-order Legendre product in $\Ihat q_k$.

\begin{proposition}[Factorization of the hybrid Vandermonde matrix]
\label{prop:factor}
$|\det(\matr V_T)|=|\det(\matr V_{\mathrm{ref}})|\,|\det(\matr S)|$. Hence, the nodes are unisolvent for $\hat\Tp_p$ if and only if $\matr S$ is regular, and they are unisolvent for every element map that is sufficiently close to an affine one.
\end{proposition}
\begin{proof}
Every row of $\matr Q$ is a vector of nodal values and has a unique expansion in the rows of the regular matrix $\matr V_{\mathrm{ref}}$, whose coefficients are the Legendre coefficients of $\Ihat q_k$, i.e., $\matr Q=\matr S^\T\matr L+\matr B\matr R$. Hence, $\matr V_T=\bigl[\begin{smallmatrix}\matr S^\T&\matr B\\\matr0&\matr I\end{smallmatrix}\bigr]\bigl[\begin{smallmatrix}\matr L\\\matr R\end{smallmatrix}\bigr]$, and the determinant identity follows. On an affine element, $q\circ\bx_K\in\Pp_m(\xi,\eta)$, $\Ihat$ is exact, and $\matr S$ is the matrix of the invertible composition with an affine map. Moreover, $\det(\matr S)$ depends continuously on the nodes.
\end{proof}

Since $\matr V_{\mathrm{ref}}$ is fixed and well conditioned, $\mathrm{rcond}(\matr S)$ is a cheap element-level indicator of unisolvence, computed from an $n_m\times n_m$ instead of an $n\times n$ matrix. On strongly non-affine (folded) elements, however, unisolvence is not guaranteed. In \cref{sec:num-cylinder}, we show that it fails on some folded elements for $p\ge7$, exactly where $\matr S$ becomes singular.
%
\subsection{Approximation}
\label{sec:hybrid-approximation}
%
The best-fit term in \cref{eq:quasiopt} is bounded by the error of the interpolant, which is a trial field with the exact nodal values. For the interpolant, the following estimate holds.

\begin{lemma}[Interpolation error]
\label{lem:interp}
Let $\mathcal U_K\in\{\Tp_p,\Qsk\}$ be unisolvent, $k_K$ its complete degree ($m$ or $p$, respectively), $\Lam_{\mathcal U}$ its Lebesgue constant, $h_K$ the diameter of $K$ and $K^\ast$ the convex hull of $K$. Then, for $f\in W^{k_K+1,\infty}(K^\ast)$,
%
\begin{equation}
  \begin{aligned}
  \|f-\mathcal I_{\mathcal U}f\|_{L^\infty(K)}
  &\le(1+\Lam_{\mathcal U})\inf_{v\in\Pp_{k_K}(x,y)}\|f-v\|_{L^\infty(K)}\\
  &\le C\,(1+\Lam_{\mathcal U})\,h_K^{k_K+1}|f|_{W^{k_K+1,\infty}(K^\ast)} .
  \end{aligned}
  \label{eq:interp}
\end{equation}
%
Moreover, with the derivative Lebesgue constant
%
\begin{equation}
  \Lam'_{\mathcal U}=h_K\max_{\bx\in K}\sum_i\bigl|\nabla M_i^K(\bx)\bigr| ,
  \label{eq:lebesgue1}
\end{equation}
%
we obtain
%
\begin{equation}
  |f-\mathcal I_{\mathcal U}f|_{W^{1,\infty}(K)}
  \le C\,(1+\Lam'_{\mathcal U})\,h_K^{k_K}|f|_{W^{k_K+1,\infty}(K^\ast)} .
  \label{eq:interp1}
\end{equation}
%
For a family of elements obtained from a fixed element by similarity transformations (self-similar refinement), $\Lam_{\mathcal U}$ and $\Lam'_{\mathcal U}$ do not depend on $h_K$.
\end{lemma}
\begin{proof}
$\mathcal I_{\mathcal U}$ reproduces $\Pp_{k_K}(x,y)$, so that $f-\mathcal I_{\mathcal U}f=(f-v)-\mathcal I_{\mathcal U}(f-v)$ for every such $v$. We choose $v$ as the Taylor polynomial of degree $k_K$ of $f$ about a point of $K^\ast$. Then, $\|f-v\|_{L^\infty}\le Ch_K^{k_K+1}|f|_{W^{k_K+1,\infty}}$ and $|f-v|_{W^{1,\infty}}\le Ch_K^{k_K}|f|_{W^{k_K+1,\infty}}$ on $K^\ast$, with $C$ depending only on $k_K$. The $L^\infty$ estimate follows from $\|\mathcal I_{\mathcal U}(f-v)\|_\infty\le\Lam_{\mathcal U}\|f-v\|_\infty$ and the $W^{1,\infty}$ estimate from $|\nabla\mathcal I_{\mathcal U}(f-v)|\le\|f-v\|_\infty\sum_i|\nabla M_i^K|\le h_K^{-1}\Lam'_{\mathcal U}\|f-v\|_\infty$. Under a similarity $\bx\mapsto\alpha\matr R\bx+\vect c$, the shape functions are unchanged as functions of the material point (\cref{lem:affine}), their gradients scale with $\alpha^{-1}$ and $h_K$ with $\alpha$, so that $\Lam'_{\mathcal U}$ is invariant.
\end{proof}

Thus, both metric spaces have the optimal order of interpolation in $L^\infty$ and in $W^{1,\infty}$, and hence in the energy norm, on meshes whose distortion persists under refinement, on which the conventional element is limited to $\Pp_{\lfloor p/q\rfloor}(x,y)$ (\cref{sec:completeness}). The two metric spaces differ in the constants $1+\Lam_{\mathcal U}$ and $1+\Lam'_{\mathcal U}$. Within \cref{eq:quasiopt}, these constants enter the best-fit term only, while the stability factor $\beta_h$ depends on the pair of test and trial spaces and is not covered by \cref{lem:interp}. An error estimate for the discrete solution would therefore require a discrete inf-sup condition for the pair of the conforming test space and the non-conforming trial space together with a consistency estimate in the sense of Strang. We do not prove such an estimate for either space. Instead, the stability is assessed numerically in \cref{sec:num-convergence,sec:num-spectra,sec:num-cylinder}.
%
\subsection{Gradients, stiffness matrix and folded elements}
\label{sec:hybrid-gradients}
%
For $\phi=L_a(\tilde\xi)L_b(\tilde\eta)$ with $a+b\le m$, the gradient is given by $\nabla_{\bx}\phi=\matr J_0^{-1}\operatorname{diag}(r_{\xi}^{-1},r_{\eta}^{-1})\,\nabla_{\bxit}\phi$. For $\phi\in\Cp_p$, on the other hand, $\nabla_{\bx}\phi=\matr J(\bxi)^{-1}\nabla_{\bxi}\phi$ at the quadrature point $\bxi$. Since the hybrid functions are evaluated through $\bxi$, the stiffness matrix \cref{eq:Ke} requires no inverse map, and post-processing at arbitrary physical points requires $\bx_K^{-1}$ only as for the conventional element. The complement contributes rational integrands to $\matr K_K$, as for curved isoparametric elements. This does not affect the exact reproduction of polynomial solutions, since for $\vect u\in[\Pp_m(x,y)]^2$ the coefficients of the interpolant on $\Cp_p$ vanish (direct sum). Then, $\matr B_{\mathrm M}\vect d=\bm\varepsilon(\vect u)$, and the integrand $\bar{\matr B}_{\mathrm N}^\T\matr C\,\bm\varepsilon(\vect u)$ is the same polynomial as for $\Qsk$, so that condition (C3) can be met by Gauss quadrature. Note that the sign of $\det(\matr J)$ is irrelevant for the function values. A vanishing $\det(\matr J)$ at a quadrature point, however, has to be excluded, since the gradients of the complement functions involve $\matr J^{-1}$. On a folded element, $\det(\matr J)$ changes sign along a curve inside the element, and the complement part of the integrand behaves like $1/\det(\matr J)$ near this curve. This singularity is not integrable. Consequently, the stiffness matrix of the hybrid space, like that of the conventional element, does not exist on folded elements, and its Gauss approximations do not converge (\cref{sec:num-quadrature}). Only the product of the stiffness matrix with the nodal values of a polynomial field of degree $m$ remains well defined, since the complement does not contribute to it.

\section{Implementation}
\label{sec:implementation}
%
Both improvements are local to the routine that evaluates the metric shape functions at a set of query points. Assembly, loads, boundary conditions and solvers remain unchanged. \Cref{alg:metric} summarizes the evaluation for one element. Here, the query points are passed with their reference coordinates, which are known to the element routines, e.g., as quadrature or plotting points.

\begin{algorithm}[htbp]
\caption{Metric shape functions of one element at the query points $\bxi_r$ (\PM: trial space $\Qsk$, Lagrange basis; \PH: hybrid trial space $\Tp_p$).}
\label{alg:metric}
\begin{algorithmic}[1]
\Require nodes $\bx_j$, reference nodes $\bxi_j$, query points $\bxi_r$, $\bx_r=\bx_K(\bxi_r)$
\State $\bx_0\gets\bx_K(\vect 0)$, $\matr J_0\gets\matr J(\vect 0)$; $\bxib_j\gets\matr J_0^{-\T}(\bx_j-\bx_0)$, $\bxib_r\gets\matr J_0^{-\T}(\bx_r-\bx_0)$
  \Comment{skew coordinates \cref{eq:skew}}
\State scale $\bxib_j,\bxib_r$ with the bounding box of $\{\bxib_j\}$ \Comment{\cref{eq:bbox}}
\If{\PM}
  \State $V_{cj}\gets\ell_a(\tilde\xi_j)\ell_b(\tilde\eta_j)$, $P_{cr}\gets\ell_a(\tilde\xi_r)\ell_b(\tilde\eta_r)$, $0\le a,b\le p$
    \Comment{1D GLL Lagrange polynomials}
  \State gradients: $\nabla_{\bx}=\matr J_0^{-1}\operatorname{diag}(r_{\xi}^{-1},r_{\eta}^{-1})\nabla_{\bxit}$
\Else \Comment{\PH, \cref{eq:hybrid}}
  \For{$0\le a\le p$, $0\le b\le p$}
    \If{$a+b\le p$}
      \State $V_{cj}\gets L_a(\tilde\xi_j)L_b(\tilde\eta_j)$, $P_{cr}\gets L_a(\tilde\xi_r)L_b(\tilde\eta_r)$; gradient as in line 5
    \Else
      \State $V_{cj}\gets L_a(\xi_j)L_b(\eta_j)$, $P_{cr}\gets L_a(\xi_r)L_b(\eta_r)$; $\nabla_{\bx}=\matr J(\bxi_r)^{-1}\nabla_{\bxi}$
    \EndIf
  \EndFor
\EndIf
\State stop if $\mathrm{rcond}(\matr V)<\varepsilon_{\mathrm{mach}}$ \Comment{shape functions do not exist}
\State LU factorization of $\matr V$; $\matr N_{\mathrm M}\gets\matr V^{-1}\matr P$, $\nabla\matr N_{\mathrm M}\gets\matr V^{-1}\nabla\matr P$
\end{algorithmic}
\end{algorithm}

The one-dimensional Lagrange basis costs $\mathcal O(p)$ operations per function and point (product formula, with prefix and suffix products for the derivative, so that query points may coincide with nodes), compared with $\mathcal O(1)$ for monomials or recurrences. The hybrid trial functions additionally need $\matr J(\bxi_r)$, which is computed by the element routines anyway. The dominant cost, namely the LU factorization of the $n\times n$ matrix $\matr V$ and the solves with the $\mathcal O(p^2)$ right-hand sides of $\matr P$ and its gradients, is of the order of $p^6$ for all variants. Measured run times are given in \cref{sec:num-cost}.

If $\matr V$ is singular to working precision, the element is rejected. The companion paper, in contrast, moves a corner node of such an element by a small fraction of its shortest adjacent edge~\cite{Eisentraeger2026axisymmetric}. Since this changes the mesh, we do not use it here and report rejected cases instead (\cref{sec:num-cylinder}).

For anisotropic degrees $p_\xi\ne p_\eta$, the trial space of \PM, $\Qp_{p_\xi,p_\eta}(\bar\xi,\bar\eta)$, contains the complete polynomials of degree $m=\min(p_\xi,p_\eta)$. The hybrid space is constructed in the same way with $m$ instead of $p$ in the first group of \cref{eq:hybrid} and with the complement $L_a(\xi)L_b(\eta)$, $a\le p_\xi$, $b\le p_\eta$, $a+b>m$. Hence, both trial spaces have the same complete degree $m$.

\section{numerics}
\begin{figure}[htbp]\centering\begin{tikzpicture}
\begin{groupplot}[group style={group size=3 by 1, horizontal sep=1.2cm},
    width=0.33\linewidth, height=0.33\linewidth, axis equal image,
    tick label style={font=\scriptsize}, title style={font=\small}]
  \nextgroupplot[title={physical $(x,y)$}, xlabel={$x$}, ylabel={$y$}]
    \addplot[gray, thick] table[x=x, y=y]{\data/outline-general.dat};
    \addplot[only marks, mark=*, mark size=1pt, blue] table[x=x, y=y]{\data/nodes-general-p8.dat};
  \nextgroupplot[title={scaled skew $(\tilde s,\tilde t)$}, xlabel={$\tilde s$}, ylabel={$\tilde t$},
      xmin=-1.15, xmax=1.15, ymin=-1.15, ymax=1.15]
    \addplot[black, dashed] coordinates {(-1,-1) (1,-1) (1,1) (-1,1) (-1,-1)};
    \addplot[gray, thick] table[x=s, y=t]{\data/outline-general.dat};
    \addplot[only marks, mark=*, mark size=1pt, red] table[x=s, y=t]{\data/nodes-general-p8.dat};
  \nextgroupplot[title={reference $(\xi,\eta)$}, xlabel={$\xi$}, ylabel={$\eta$},
      xmin=-1.15, xmax=1.15, ymin=-1.15, ymax=1.15]
    \addplot[gray, thick] coordinates {(-1,-1) (1,-1) (1,1) (-1,1) (-1,-1)};
    \addplot[only marks, mark=*, mark size=1pt, blue] table[x=xi, y=eta]{\data/nodes-general-p8.dat};
\end{groupplot}
\end{tikzpicture}
\caption{nodes}\end{figure}
\begin{figure}[htbp]\centering\begin{tikzpicture}
\begin{groupplot}[group style={group size=2 by 2, horizontal sep=1.5cm, vertical sep=1.6cm},
    width=0.47\linewidth, height=0.34\linewidth, ymode=log, grid=major,
    xlabel={$p$}, ylabel={$\kappa_2(\matr V)$}, xmin=1, xmax=16, ymin=0.5, ymax=1e17,
    cycle list name=bases, every axis plot/.append style={mark size=1.3pt},
    unbounded coords=discard]
  \nextgroupplot[title={parallelogram}, legend to name=leg:kappa, legend columns=5]
    \pgfplotsforeachungrouped \b in {mon,leg,che,lag,hyb} {\addplot table[x=p, y=kappa2_\b]{\data/element-parallelogram.dat};}
    \legend{monomial, Legendre, Chebyshev, Lagrange (\PM), hybrid (\PH)}
  \nextgroupplot[title={trapezoid}]
    \pgfplotsforeachungrouped \b in {mon,leg,che,lag,hyb} {\addplot table[x=p, y=kappa2_\b]{\data/element-trapezoid.dat};}
  \nextgroupplot[title={curved}]
    \pgfplotsforeachungrouped \b in {mon,leg,che,lag,hyb} {\addplot table[x=p, y=kappa2_\b]{\data/element-curved.dat};}
  \nextgroupplot[title={general}]
    \pgfplotsforeachungrouped \b in {mon,leg,che,lag,hyb} {\addplot table[x=p, y=kappa2_\b]{\data/element-general.dat};}
\end{groupplot}
\end{tikzpicture}\\[3pt]
\pgfplotslegendfromname{leg:kappa}
\caption{kappa}\end{figure}
\begin{figure}[htbp]\centering\begin{tikzpicture}
\begin{groupplot}[group style={group size=2 by 2, horizontal sep=1.5cm, vertical sep=1.6cm},
    width=0.44\linewidth, height=0.34\linewidth, ymode=log, grid=major,
    xlabel={$p$}, ylabel={$e_{\mathrm{K}}$}, xmin=1, xmax=16, ymin=1e-17, ymax=1,
    cycle list name=bases, every axis plot/.append style={mark size=1.3pt},
    unbounded coords=discard, y filter/.expression={y<-39.2 ? nan : y}]
  \nextgroupplot[title={parallelogram}]
    \pgfplotsforeachungrouped \b in {mon,leg,che,lag,hyb} {\addplot table[x=p, y=eKron_\b]{\data/element-parallelogram.dat};}
  \nextgroupplot[title={trapezoid}]
    \pgfplotsforeachungrouped \b in {mon,leg,che,lag,hyb} {\addplot table[x=p, y=eKron_\b]{\data/element-trapezoid.dat};}
  \nextgroupplot[title={curved}]
    \pgfplotsforeachungrouped \b in {mon,leg,che,lag,hyb} {\addplot table[x=p, y=eKron_\b]{\data/element-curved.dat};}
  \nextgroupplot[title={general}]
    \pgfplotsforeachungrouped \b in {mon,leg,che,lag,hyb} {\addplot table[x=p, y=eKron_\b]{\data/element-general.dat};}
\end{groupplot}
\end{tikzpicture}\\[3pt]
\pgfplotslegendfromname{leg:kappa}
\caption{kron}\end{figure}
\begin{figure}[htbp]\centering\begin{tikzpicture}
\begin{groupplot}[group style={group size=3 by 1, horizontal sep=1.1cm},
    width=0.36\linewidth, height=0.32\linewidth, ymode=log, grid=major,
    xlabel={$p$}, xmin=1, xmax=16, ymin=1, ymax=1e10, unbounded coords=discard]
  \nextgroupplot[title={trapezoid}, ylabel={Lebesgue constant},
      legend to name=leg:leb, legend columns=3]
    \addplot[styPM] table[x=p, y=LQ]{\data/element-trapezoid.dat};
    \addplot[styPH] table[x=p, y=LT]{\data/element-trapezoid.dat};
    \addplot[styLN] table[x=p, y=LN]{\data/element-trapezoid.dat};
    \legend{{$\Lam_Q$, $\Qp_p(s,t)$}, {$\Lam_T$, $\Tp_p$}, {$\Lam_N$, GLL}}
  \nextgroupplot[title={curved}]
    \addplot[styPM] table[x=p, y=LQ]{\data/element-curved.dat};
    \addplot[styPH] table[x=p, y=LT]{\data/element-curved.dat};
    \addplot[styLN] table[x=p, y=LN]{\data/element-curved.dat};
  \nextgroupplot[title={general}]
    \addplot[styPM] table[x=p, y=LQ]{\data/element-general.dat};
    \addplot[styPH] table[x=p, y=LT]{\data/element-general.dat};
    \addplot[styLN] table[x=p, y=LN]{\data/element-general.dat};
\end{groupplot}
\end{tikzpicture}\\[3pt]
\pgfplotslegendfromname{leg:leb}
\caption{leb}\end{figure}

\section{Discussion}
\label{sec:discussion}

\paragraph{Effect of the two improvements}
Note that the two improvements act on different sources of error and, therefore, have to
be assessed separately. The Lagrange basis of \cref{sec:basis} removes the
ill-conditioning \emph{of the construction}. It is exact on affine elements, and on
distorted ones it reduces the condition number by orders of magnitude relative to
monomials. The remainder stems from extrapolating the metric interpolant to the bounding
box of the element (\cref{prop:kappa}). The basis cannot change the shape functions
(\cref{prop:invariance}), and consequently it cannot change the Lebesgue constant, the
patch-test behaviour in exact arithmetic, the eigenvalues or the discrete solution beyond
round-off. The hybrid space of \cref{sec:hybrid}, in contrast, changes the \emph{space},
and with it the Lebesgue constant, the stability properties and the approximation
constant. Both improvements are needed, since the hybrid space has its own Vandermonde
matrix. In the Legendre-type basis used here, this matrix is better conditioned than that
of $\Qp_p(s,t)$ in any basis for $p\ge4$ (by up to five orders of magnitude at $p=16$),
because the complement is evaluated on the exact GLL tensor grid. Its condition number
still grows with $p$ ($\approx2\cdot10^{6}$ at $p=16$, \cref{fig:kappa}).

\paragraph{Trade-off between the two trial spaces}
For smooth solutions on meshes where $\Qp_p(s,t)$ is stable, \PM\ is
more accurate than \PH\ by a factor of $1.3$--$5.5$ in the energy norm, at the same
convergence order (\cref{sec:num-convergence}). This is due to the larger Cartesian
polynomial content of $\Qp_p(s,t)$. This advantage is, however, lost abruptly when the
Lebesgue constant of $\Qp_p(s,t)$ becomes large. On the $\delta=0.2$ meshes, \PM\ diverges
for $p\ge5$, and on all three $3\times2$ meshes it develops complex and negative
eigenvalues from $p=5$--$7$ on. Neither the degree nor the amount of distortion at which
this happens is known a priori, and the Jacobian is not an indicator either (the bent
element has $\det\matr J\equiv1$). \PH\ was stable and optimally convergent in every
non-folded case studied. In our opinion, this makes $\Tp_p$ the appropriate default for
high-order unsymmetric spectral elements, in particular for transient, modal and
frequency-domain analyses. The space $\Qp_p(s,t)$, on the other hand, remains attractive
for static analyses at low to moderate degree and distortion, where its larger Cartesian
content pays off. An element-level indicator for the latter choice is $\Lam_Q$ itself,
which can be estimated by evaluating the metric shape functions on a sampling grid in the
element. Note that the condition number of $\matr V$ is not a substitute, since it also
contains the extrapolation to the bounding box (\cref{prop:kappa}).

\paragraph{Folded elements}
Exactness for polynomial solutions is observed on folded elements for both
spaces (\cref{rem:folded}), as the patch tests on the panel meshes of
\cref{sec:num-patch} confirm. At $p=8$, \PH\ reduces the round-off level there by one
to two orders of magnitude. For non-polynomial solutions on folded meshes, however, \PH\
converges less regularly than on valid meshes, and its energy error stagnates. This is due
to the fact that the complement part of $\matr K_e$ contains $1/\det\matr J$, which is
singular on the curve $\det\matr J=0$ inside the element. A remedy would be to integrate
these terms with a rule adapted to the singular curve or to write the complement
gradients in a form that avoids $\matr J^{-1}$ (e.g.\ by multiplying the complement
functions by $\det\matr J$, which changes the space). This is out of the scope of this
contribution and left for future research. On folded meshes, \PM\ is unstable at
high order in our tests, and \PP\ converges with reduced order in the nodal error.

\paragraph{Limitations}
The following limitations of the present study have to be kept in mind:
\begin{enumerate}
\item No a priori stability estimate (discrete inf-sup condition
  \citep{Babuska1971Error}) is available for either metric space. The quasi-optimality of
  \PH\ is observed for the thick cylinder (\cref{tab:quasiopt}), but not proven, and the
  statements on stability rest on the computed spectra, condition numbers and convergence
  histories. While \cref{prop:representation} gives exact representations and upper
  bounds of the Lebesgue constants, the size of the correction operators is only
  quantified numerically.
\item The completeness degree of $\Tp_p$ is $\min(p_\xi,p_\eta)$, the same as for
  $\Qp_{p_\xi,p_\eta}(s,t)$, and neither space reproduces
  $\Pp_{\max(p_\xi,p_\eta)}$.
\item The complement contributes rational integrands to $\matr K_e$ and
  $\matr M_e$. As for curved isoparametric elements, Gauss quadrature is then not exact.
  Where $\det\matr J$ vanishes inside an element, the integrands are singular.
\item The study is restricted to linear plane elasticity on quadrilaterals with at most
  biquadratic geometry. The construction carries over to hexahedra
  ($\Pp_p(x,y,z)\oplus\Cp_p$ with the $L^2$-complement of $\Pp_p(\xi,\eta,\zeta)$ in
  $\Qp_p(\xi,\eta,\zeta)$) and to triangles (complement in a simplex basis), but it has not
  been tested there.
\end{enumerate}

\section{Conclusions}
\label{sec:conclusions}

Two improvements of the metric trial functions of unsymmetric spectral elements were
presented, analysed and tested.
\begin{enumerate}
\item \emph{Basis of the metric Vandermonde matrix.} The metric shape functions are
  basis independent. The Lagrange polynomials through the reference GLL nodes, in
  bounding-box scaled skew coordinates, give $\cond(\matr V)=1$ on affine elements and a
  condition number equal to the Lebesgue function of the metric interpolation (up to a
  bounded factor) on distorted ones. Compared with monomials this reduces the error of the
  computed shape functions by two to four orders of magnitude on distorted elements at
  $p=16$ and restores the exact symmetry of $\matr K_e$ on affine elements. The Lebesgue
  constant, and with it everything that depends only on the space, is unaffected by
  any choice of basis.
\item \emph{Hybrid trial space.} Writing only $\Pp_p(x,y)$ in metric coordinates and the
  canonical ($L^2$-orthogonal) complement in parametric coordinates retains completeness,
  affine invariance and distortion immunity, reduces to the classical element on
  straight-edged quadrilaterals, and keeps the Lebesgue constant at or close to the GLL
  value for moderate curved-edge distortion. In the tests it removed the instabilities of
  the $\Qp_p(s,t)$ element---divergence under mesh refinement in statics, negative and
  strongly complex eigenvalues, exponential growth in time---while converging with the
  optimal order. It trades this robustness for a moderately larger error constant on
  smooth static problems.
\end{enumerate}
Open questions are a stability theory (discrete inf-sup condition) for both metric spaces,
an accurate treatment of the singular integrands of the hybrid space on folded elements,
the residual (slow) growth of the hybrid element in long-time transient analysis, and the
extension to three dimensions, triangles and nonlinear problems.

%
% CHECK: acknowledge funding, if any, e.g.:
%\paragraph{Acknowledgment} The authors gratefully acknowledge the financial support provided by the German Research Foundation (Deutsche Forschungsgemeinschaft, DFG) through the research grant EI 1188/3-1 (project number: 497531141).
%
% -----------------------------------------------------------------------------
% Statements required by many journals at submission (uncomment and adapt)
% -----------------------------------------------------------------------------
%
% CRediT author contributions (https://credit.niso.org); use only these 14 roles:
% Conceptualization, Data curation, Formal analysis, Funding acquisition,
% Investigation, Methodology, Project administration, Resources, Software,
% Supervision, Validation, Visualization, Writing -- original draft,
% Writing -- review & editing
%\paragraph{CRediT authorship contribution statement}
%\textbf{Sascha Eisentr\"ager:} Conceptualization, Methodology, Software, Formal analysis, Investigation, Writing -- original draft, Writing -- review \& editing.
%
%\paragraph{Declaration of competing interest} The authors declare that they have no known competing financial interests or personal relationships that could have appeared to influence the work reported in this paper.
%
\paragraph{Data availability} The \texttt{MATLAB}/GNU Octave code, the meshes, the scripts of all numerical studies and the raw data required to reproduce the results of this paper will be archived on Zenodo and made available in the public repository \url{https://github.com/SaschaEisentraeger/USEM-MATLAB} upon acceptance of the paper. The code was developed with the assistance of \texttt{Claude} (Anthropic) via \texttt{Claude Code}. The scripts that produce every number, table and figure of this paper are contained in the folder \texttt{manuscript/studies} of the repository, the raw data in \texttt{manuscript/data} and \texttt{studies/thick\_cylinder/results/raw\_calibrated}, and the source of every table is noted in the LaTeX source.
% CHECK: archive this state of the repository on Zenodo, make the GitHub repository public and add the Zenodo DOI here (author, 30 Oct 2026: same procedure as for the companion paper)
%
\paragraph{Declaration of generative AI and AI-assisted technologies in the manuscript preparation process} During the preparation of this work, the author used \texttt{Claude} (Anthropic), accessed through \texttt{Claude Code}, as an AI-assisted tool and retained full control over its use. The research ideas, methodology and numerical studies were defined by the author, who specified each task and reviewed, tested and approved all changes before their adoption. The tool was used to extend and document the \texttt{MATLAB}/GNU Octave code, to execute the numerical studies and post-process the results, to generate figures from computed data with reproducible scripts, to draft and revise the manuscript text in accordance with the author's instructions and to check the cited literature. After using this tool, the author reviewed and edited the content as needed and takes full responsibility for the content of the published article.
% CHECK: adapted from the companion paper; confirm that it describes how the tool was used for this paper (and switch to "authors" once co-authors are added)
%
% Equations, figures and tables are numbered per appendix section: (A.1),
% Figure A.1, Table A.1.
\appendix
\numberwithin{equation}{section}
\counterwithin{figure}{section}
\counterwithin{table}{section}
%
\section{Test elements and meshes}
\label{app:elements}
%
The single test elements are defined by 4- or 8-node input elements (corner nodes counter-clockwise, followed by the mid-side nodes of the edges 1--2, 2--3, 3--4 and 4--1) and lifted to degree $p$ by isoparametric interpolation of the input geometry:
\begin{center}
  \small
  \begin{tabular}{@{}ll@{}}
    \toprule
    element & nodes $(x,y)$ \\
    \midrule
    parallelogram & $(0,0)$, $(2,0.3)$, $(2.6,1.5)$, $(0.6,1.2)$ \\
    trapezoid & $(0,0)$, $(2,0)$, $(1.4,1)$, $(0.2,1.3)$ \\
    curved & $(0,0)$, $(2,0)$, $(2,1)$, $(0,1)$; $(1,-0.2)$, $(2.25,0.5)$, $(1,1)$, $(0,0.5)$ \\
    general & $(0,0)$, $(2,0.1)$, $(1.5,1.2)$, $(0.2,1)$; $(1.1,-0.15)$, $(1.9,0.6)$, $(0.8,1.25)$, $(0.05,0.5)$ \\
    \bottomrule
  \end{tabular}
\end{center}
The meshes of \cref{sec:num-convergence} consist of $N\times N$ 8-node elements on the unit square whose interior mid-side nodes are moved normal to their edge by $\delta h$ with $h=1/N$, so that the relative sagitta of every interior edge is $\delta$. The normal is oriented consistently, so that all interior horizontal edges bow in the $-y$-direction and all interior vertical edges in the $-x$-direction. The $3\times2$ meshes of \cref{sec:num-patch,sec:num-spectra} on $[0,3]\times[0,2]$ use the same curved-edge distortion with $\delta=0.15$ and $0.20$ (with $h=1$). For the mesh denoted by \emph{angular + curved $0.30$}, the interior corner nodes are first shifted alternately by $\pm0.3$ in the $x$-direction, before the curved-edge distortion with $\delta=0.15$ is applied.
%
\section{Derivations}
\label{app:derivations}
%
For the sake of completeness, this appendix collects the derivations of the statements used in \cref{sec:formulation,sec:basis,sec:hybrid}. All functions of an element are considered as functions on the reference square $\Kh$, composed with the geometry map $\bx_K$. We write $\Qp=\Qp_p(\xi,\eta)$, $\Pp_p(\xi,\eta)$ for the polynomials of total degree $p$ in the reference coordinates, $\hat\Pp_p=\{q\circ\bx_K:q\in\Pp_p(x,y)\}$ for the complete polynomials of degree $p$ in $(x,y)$ and $\hat\Qp=\{q\circ\bx_K:q\in\Qsk\}$ for the trial space of \PM. $\Ihat$ denotes the GLL interpolation onto $\Qp$, and $\mathcal I_{\mathrm{T}}$ and $\mathcal I_{\mathrm{Q}}$ denote the nodal interpolation by the trial functions of \PH\ and \PM. For anisotropic degrees, $p$ is replaced by $m=\min(p_\xi,p_\eta)$ in $\Pp_p$ and $\hat\Pp_p$, and all arguments carry over.

\paragraph{Skew coordinates under affine transformations (\cref{sec:skew}).} If the nodes are transformed by $\bx\mapsto\matr A\bx+\vect c$ with $\det(\matr A)\ne0$, the geometry map becomes $\bx_K'=\matr A\bx_K+\vect c$. Hence, $\matr J_0'^\top=\matr A\matr J_0^\top$ and $\bx'-\bx_0'=\matr A(\bx-\bx_0)$, so that $\bxib'=\matr J_0^{-\top}\matr A^{-1}\matr A(\bx-\bx_0)=\bxib$. The bounding box \cref{eq:bbox} is computed from the skew coordinates of the nodes and is therefore unchanged as well.

\paragraph{Basis invariance (\cref{sec:basis-invariance}).} For another basis $\bm\psi=\matr G\bm\phi$ with a regular matrix $\matr G$, we obtain $\matr V_\psi=\matr G\matr V_\phi$ and $\matr P_\psi=\matr G\matr P_\phi$, so that $\matr V_\psi^{-1}\matr P_\psi=\matr V_\phi^{-1}\matr G^{-1}\matr G\matr P_\phi=\matr V_\phi^{-1}\matr P_\phi$. The functions $M_i^K=\sum_c(\matr V^{-1})_{ic}\phi_c$ lie in the trial space and satisfy $M_i^K(\bx_j)=\delta_{ij}$. If $\matr V$ is regular, they are the only such functions. For the nodal data $\vect f=\matr V^\top\vect c$ of a function of the trial space with the coefficient vector $\vect c$, the first-order error of the computed interpolant is $\vect c^\top\delta\matr V\,\matr N_{\mathrm M}$, where $\delta\matr V$ denotes the backward error of the factorization. Up to a constant, this error is bounded by $u\,\|\vect c\|\,\|\matr V\|\,\Lam_{\mathcal U}$. For polynomials of low degree, $\|\vect c\|\,\|\matr V\|$ is moderate in every basis, whereas generic nodal data can have $\|\vect c\|\approx\|\matr V^{-1}\|\,\|\vect f\|$.

\paragraph{Condition number in the Lagrange basis, \cref{eq:kappa}.} With the Lagrange basis, $\phi_c(\bm g_d)=\delta_{cd}$. Choosing the query points $\bm g_d$ in \cref{eq:vandermonde} gives $\matr P=\matr I$ and hence $\matr N_{\mathrm M}=\matr V^{-1}$, which gives $(\matr V^{-1})_{id}=M_i^K(\bm g_d)$. The 1-norm is the largest absolute column sum. Thus, $\|\matr V\|_1=\max_j\sum_c|\phi_c(\bxit_j)|=\max_j\lambda_{\mathrm{ref}}(\bxit_j)$ and $\|\matr V^{-1}\|_1=\max_d\sum_i|M_i^K(\bm g_d)|=\max_d\lambda_{\mathcal U}(\bm g_d)$.

\paragraph{Orthogonal complement (\cref{sec:hybrid-definition}).} Let $a+b>p$ and $i+j\le p$. Then $i<a$ or $j<b$, and the orthogonality of $L_a$ to all polynomials of degree less than $a$ gives $\int_{\Kh}L_a(\xi)L_b(\eta)\,\xi^i\eta^j\dd\xi\dd\eta=\int L_a\xi^i\dd\xi\int L_b\eta^j\dd\eta=0$. The functions $L_aL_b$ with $a,b\le p$ are linearly independent and span $\Qp$, those with $a+b\le p$ span $\Pp_p(\xi,\eta)$, and $\dim\Cp_p=n-\dim\Pp_p$. Hence, $\Cp_p$ is the $L^2(\Kh)$-orthogonal complement of $\Pp_p(\xi,\eta)$ in $\Qp$, which is defined independently of any basis.

\paragraph{Properties of the hybrid trial space (\cref{sec:hybrid-properties}).} (i) The hybrid space $\hat\Tp_p=\hat\Pp_p+\Cp_p$ is spanned by $n$ functions whose values at the $n$ nodes form the regular matrix $\matr V_{\mathrm{T}}$. Hence, they are linearly independent, the sum is direct, and the nodal map is bijective. (ii) and (iii) follow immediately from the definition and the invariance of the skew coordinates and of $\bxi$. (iv) For a bilinear map, $x,y\in\Qp_1(\xi,\eta)$, and hence $x^ay^b\circ\bx_K\in\Qp_{a+b}(\xi,\eta)\subset\Qp$ for $a+b\le p$. With $\Cp_p\subset\Qp$, we obtain $\hat\Tp_p\subset\Qp$, and equality follows from (i).

\paragraph{Representation of the interpolants, \cref{eq:repT}.} Let $\pi$ denote the projector of $\Qp$ onto $\Pp_p(\xi,\eta)$ along $\Cp_p$, i.e., the truncation of the Legendre expansion to $a+b\le p$, and define $S=\pi\Ihat|_{\hat\Pp_p}:\hat\Pp_p\to\Pp_p(\xi,\eta)$ and $S_{\mathrm{Q}}=\Ihat|_{\hat\Qp}:\hat\Qp\to\Qp$. We assume that $\bx_K$ has a point with $\det(\matr J)\ne0$, so that $\dim\hat\Pp_p=\dim\Pp_p$. The nodal values of a function $\hat v$ are in one-to-one correspondence with $\Ihat\hat v\in\Qp$, and $\Ihat c=c$ for $c\in\Cp_p\subset\Qp$. Let $\hat v=q+c$ with $q\in\hat\Pp_p$ and $c\in\Cp_p$ have the nodal values of $f$, i.e., $\Ihat q+c=\Ihat f$. Applying $\pi$, whose kernel is $\Cp_p$, gives $Sq=\pi\Ihat f$. Conversely, for every $q$ with $Sq=\pi\Ihat f$, the function $c=\Ihat f-\Ihat q$ lies in $\ker\pi=\Cp_p$. If $S$ is invertible, $q$ and $c$ are therefore uniquely determined for every $f$, and $\hat v=q+\Ihat f-\Ihat q$, which is \cref{eq:repT} with $q=S^{-1}\pi\Ihat f$. If $S$ is singular, we take $0\ne q\in\ker S$ and $c=-\Ihat q\in\Cp_p$. The function $\hat v=q+c$ then vanishes at all nodes. Either $\hat v\ne0$, or $q=-c\in\hat\Pp_p\cap\Cp_p$, so that the sum is not direct and $\dim\hat\Tp_p<n$. In both cases, the hybrid shape functions do not exist. Thus, the hybrid space is unisolvent at the nodes if and only if $S$ is invertible. Likewise, an interpolant $\hat v\in\hat\Qp$ satisfies $\Ihat\hat v=\Ihat f$, i.e., $S_{\mathrm{Q}}\hat v=\Ihat f$, and $\hat v=\Ihat\hat v+(I-\Ihat)\hat v$, so that $\mathcal I_Qf=\Ihat f+(I-\Ihat)S_{\mathrm{Q}}^{-1}\Ihat f$. With $\|\Ihat f\|_\infty\le\Lam_{\mathrm{N}}\|f\|_\infty$, we obtain the bounds
%
\begin{equation}
  \Lam_{\mathrm{T}}\le\Lam_{\mathrm{N}}\bigl(1+\|(I-\Ihat)S^{-1}\pi\|\bigr)\qquad\text{and}\qquad
  \Lam_{\mathrm{Q}}\le\Lam_{\mathrm{N}}\bigl(1+\|(I-\Ihat)S_{\mathrm{Q}}^{-1}\|\bigr),
  \label{eq:lebbound}
\end{equation}
%
with operator norms on $\Qp$ in the maximum norm on $\Kh$. For the nodal data of node $i$, $\Ihat f=\hat N_i$, so that the shape functions are $\hat M_{i,\mathrm{T}}=\hat N_i+(I-\Ihat)S^{-1}\pi\hat N_i$ and $\hat M_{i,\mathrm{Q}}=\hat N_i+(I-\Ihat)S_{\mathrm{Q}}^{-1}\hat N_i$. If $\hat\Pp_p\subset\Qp$, in particular for every bilinear map, $(I-\Ihat)q=0$, $\mathcal I_{\mathrm{T}}=\Ihat$ and $\Lam_{\mathrm{T}}=\Lam_{\mathrm{N}}$.

\paragraph{Factorization of the hybrid Vandermonde matrix, \cref{eq:factor}.} Let the rows of $\matr V_{\mathrm{ref}}$ be the nodal values of $L_a(\xi)L_b(\eta)$, $a,b\le p$, ordered into the $n_p$ low-order rows $\matr L$ ($a+b\le p$) and the rows $\matr R$ of the complement. The hybrid Vandermonde matrix is ordered alike, $\matr V_{\mathrm{T}}=[\matr Q;\matr R]$, where the rows of $\matr Q$ are the nodal values of a basis $\{q_k\}$ of $\hat\Pp_p$. Every row of $\matr Q$ has a unique expansion in the rows of the regular matrix $\matr V_{\mathrm{ref}}$, whose coefficients are the Legendre coefficients of $\Ihat q_k$, i.e., $\matr Q=\matr S^\top\matr L+\matr B\matr R$, where $\matr S$ is the matrix of $S$ and $\matr B$ contains the coefficients of the complement. Hence, $\matr V_{\mathrm{T}}=\bigl[\begin{smallmatrix}\matr S^\top&\matr B\\\matr0&\matr I\end{smallmatrix}\bigr]\bigl[\begin{smallmatrix}\matr L\\\matr R\end{smallmatrix}\bigr]$, and \cref{eq:factor} follows. On an affine element, $q\circ\bx_K\in\Pp_p(\xi,\eta)$, $\Ihat$ is exact, and $\matr S$ is the matrix of the invertible composition with an affine map. Since $\det(\matr S)$ depends continuously on the nodes, $\matr S$ is regular on every element that is sufficiently close to an affine one.

\paragraph{Interpolation estimates, \cref{eq:interp}.} Both trial spaces reproduce $\Pp_p(x,y)$, so that $f-\mathcal I_{\mathcal U}f=(f-v)-\mathcal I_{\mathcal U}(f-v)$ for every $v\in\Pp_p(x,y)$. We choose $v$ as the Taylor polynomial of degree $p$ of $f$ about a point of the convex hull $K^\ast$ of $K$. Then, $\max_{K^\ast}|f-v|\le Ch_K^{p+1}|f|_{p+1}$ and $\max_{K^\ast}|\nabla(f-v)|\le Ch_K^{p}|f|_{p+1}$, with $C$ depending only on $p$. The first estimate in \cref{eq:interp} follows from $\max_K|\mathcal I_{\mathcal U}(f-v)|\le\Lam_{\mathcal U}\max_K|f-v|$ and the second from $|\nabla\mathcal I_{\mathcal U}(f-v)|\le\max_K|f-v|\sum_i|\nabla M_i^K|\le h_K^{-1}\Lam'_{\mathcal U}\max_K|f-v|$. Under a similarity transformation $\bx\mapsto\alpha\matr R\bx+\vect c$, the shape functions are unchanged as functions of the material point, their gradients scale with $\alpha^{-1}$ and $h_K$ with $\alpha$, so that $\Lam_{\mathcal U}$ and $\Lam'_{\mathcal U}$ are invariant.

%
% -----------------------------------------------------------------------------
% References
% -----------------------------------------------------------------------------
\bibliographystyle{./bib/IEEEtranDOI}
\phantomsection
\addcontentsline{toc}{section}{References}
\bibliography{./bib/references}
%
\end{document}
